\chapter{数据工程——万亿 Token 的炼金术}
\label{ch:data}
%% ============================================================

如果说 Transformer 是 LLM 的大脑，
那么数据就是它的食粮。
一个 LLM 的质量上限在很大程度上取决于训练数据的质量和多样性。
2025--2026 年的共识越来越清晰：
\textbf{数据质量远比数据数量重要}。

这不是一句空洞的口号。
微软的 Phi 系列模型用不到 100B 的精选高质量 token
训练出了与万亿 token 模型竞争的结果，
而大量研究表明在未经过滤的网页数据上
增加训练量带来的收益迅速递减。
数据工程：从原始互联网爬取数据到
可以喂给模型的高质量 token 流：
是 LLM 训练中最不起眼却最关键的环节。

\section{数据来源}

训练一个 frontier LLM 需要万亿级别的 token。
这些数据来自哪里？
让我们逐一审视每种数据来源的特性、规模和挑战。

\subsection{网页数据：最大的宝藏与垃圾场}

\textbf{Common Crawl} 是几乎所有 LLM 训练数据的源头。
这个非营利组织自 2008 年起持续爬取互联网，
每月发布一个新的爬取快照，
累积了数百 PB 的原始 HTML 数据，
覆盖超过\textbf{2500 亿个网页}。
每个月的快照以 WARC（Web ARChive）格式存储，
包含完整的 HTTP 响应（HTML、头信息等），
同时提供 WET 格式（仅提取文本内容）
和 WAT 格式（仅元数据）。

但原始 Common Crawl 的质量极差：
充满了导航栏、广告、重复文本、
SEO 垃圾、色情内容和机器生成的胡言乱语。
直接使用原始 Common Crawl 训练模型
得到的效果远不如在清洗后的子集上训练。
因此，围绕 Common Crawl 衍生出了一系列
精心策划的高质量数据集：

\begin{itemize}[noitemsep]
  \item \textbf{C4}（Colossal Clean Crawled Corpus）：
        Google 为 T5 模型创建，
        对 Common Crawl 2019 年 4 月快照进行清洗，
        得到约 750GB / 156B token 的英文文本。
        清洗规则相对简单：
        去除含脏词的页面、过短的行、
        含 “lorem ipsum” 的页面等。
        % NEW REF: raffel2020t5 = Raffel, Colin and Shazeer, Noam and Roberts, Adam and Lee, Katherine and Narang, Sharan and Matena, Michael and Zhou, Yanqi and Li, Wei and Liu, Peter J. Exploring the Limits of Transfer Learning with a Unified Text-to-Text Transformer. JMLR, 21(140):1--67, 2020.
  \item \textbf{mC4}：C4 的多语言版本，
        覆盖 100+ 种语言。
  \item \textbf{RefinedWeb}（Falcon 团队）：
        对 Common Crawl 进行了更精细的清洗和去重，
        得到约 5T token，证明了仅用网页数据
        （不混入书籍或代码）也能训练出高质量模型。
        % NEW REF: penedo2023refinedweb = Penedo, Guilherme and Malartic, Quentin and Hesslow, Daniel and others. The RefinedWeb Dataset for Falcon LLM. In NeurIPS Datasets and Benchmarks Track, 2023.
  \item \textbf{FineWeb}~\cite{penedo2024fineweb}（HuggingFace）：
        目前最大、最高质量的公开网页数据集。
        对 96 个 Common Crawl 快照（2013--2024 年）
        进行了精细清洗，
        得到了 \textbf{15T token} 的高质量英文文本。
        FineWeb-Edu 是其教育内容子集（1.3T token），
        在多个基准测试上的训练效果
        甚至优于使用全量 FineWeb。
  \item \textbf{DCLM}（DataComp-LM，Allen AI 等）：
        通过系统性的受控实验筛选最优的过滤策略，
        从 Common Crawl 中提取了 3T+ token 的高质量数据。
        % NEW REF: li2024dclm = Li, Jeffrey and others. DataComp-LM: In search of the next generation of training sets for language models. In NeurIPS, 2024.
  \item \textbf{FineWeb-2}（HuggingFace，2024 年 12 月）：
        FineWeb 的多语言继任者。
        FineWeb-2 将 FineWeb 的英文处理流水线
        推广到 \textbf{1000+ 种语言}，
        产生了约 3T 词（压缩后约 8TB）的多语言高质量数据。
        其核心贡献在于解决了多语言数据处理中的独特挑战：
        不同语言需要不同的质量过滤阈值、
        低资源语言的语言识别准确率较低、
        某些书写系统（如阿拉伯文、天城文）
        需要特殊的预处理规则。
        FineWeb-2 的论文（2025 年 6 月）
        详细记录了这些挑战和解决方案，
        是构建多语言预训练数据的重要参考。
        % NEW REF: penedo2025fineweb2 = Penedo, Guilherme and Kydlicek, Hynek and Sabolcec, Vinko and Messmer, Bettina and Foroutan, Negar and Kargaran, Amir Hossein and Raffel, Colin and Jaggi, Martin and Von Werra, Leandro and Wolf, Thomas. FineWeb2: One Pipeline to Scale Them All. arXiv preprint arXiv:2506.20920, 2025.
  \item \textbf{Nemotron-CC}（NVIDIA，2025 年 1 月）：
        NVIDIA 构建的万亿 token 级英文预训练数据集。
        Nemotron-CC 的独特之处在于其
        \textbf{分类器集成 + 合成数据改写}的双重策略：
        首先使用多个质量分类器的集成
        （而非单个分类器）来评估文档质量，
        然后对中等质量的文档使用 LLM 进行\textbf{改写}
        （rephrasing），将其提升为高质量版本。
        实验显示，在 Nemotron-CC 上训练的 8B 参数模型
        在多项任务上超越了 LLaMA 3.1 8B。
        这种「改写提质」的策略
        为数据质量工程提供了一个新的维度：
        \textbf{不只是过滤掉低质量数据，
        还可以主动改善数据质量}。
        % NEW REF: nvidia2025nemotroncc = Su, Dan and others. Nemotron-CC: Transforming Common Crawl into a Refined Long-Horizon Pretraining Dataset. arXiv preprint arXiv:2501.05046, 2025.
\end{itemize}

\subsection{书籍：长文本的稀缺宝藏}

\textbf{书籍}提供长篇连贯的高质量文本，
尤其对模型的长文本生成和叙事能力有帮助。
书籍文本的质量通常远高于网页：
经过专业编辑和出版流程的审核。

GPT-3~\cite{brown2020gpt3} 的训练使用了
Books1 和 Books2 两个书籍数据集
（共约 67B token），
但 OpenAI 从未公开其具体来源。
The Pile（EleutherAI 的开源数据集）
包含了 Books3 子集：
约 19.7 万本书的全文，来自盗版电子书网站。
这引发了巨大的版权争议，
多位作家和出版社对此提起了诉讼，
Books3 最终被从公开发布的版本中移除。

版权问题使得合法获取大量书籍数据变得困难。
目前的替代方案包括：

\begin{itemize}[noitemsep]
  \item \textbf{Project Gutenberg}：
        约 7 万本公共领域的经典书籍
        （版权已过期的作品，主要是 19 世纪及更早的文学作品）。
        数据质量极高（经过人工校对），
        但时间分布严重偏向过去：
        模型从中学到的是 “经典英语” 而非现代语言。
  \item \textbf{Internet Archive / Open Library}：
        提供数百万本书的借阅式数字访问。
        2023 年的 Hachette v. Internet Archive 判决
        裁定其“受控数字借阅”模式
        不构成合理使用：
        这一判决进一步限制了
        通过 Internet Archive 获取书籍数据的可能性。
  \item \textbf{出版商许可合作}：
        OpenAI 与 Springer Nature、Wiley 等出版商
        签订了多年期的内容许可协议，
        据报道每年费用在数百万到数千万美元级别。
        这种“付费获取”模式正在成为
        frontier 模型训练的标准做法。
  \item \textbf{LibGen 和盗版来源}：
        尽管法律风险显而易见，
        社区中仍有讨论认为
        部分模型的训练数据可能包含来自
        盗版渠道的书籍数据：
        Books3 就是一个已被证实的例子。
        随着版权诉讼的推进，
        这种做法的法律风险在急剧上升。
\end{itemize}

书籍数据在整体训练数据中的占比通常为 5--10\%：
数量不多，但对模型质量的贡献\textbf{不成比例地高}。
这种“少量高杠杆”的特性有一个合理的解释：
书籍提供了互联网数据中极为稀缺的\textbf{长程连贯性}：
一本书的数十万词围绕一个或几个主题展开，
有严密的逻辑链条和叙事弧线。
相比之下，网页文本通常只有数百到数千词，
且不同网页之间缺乏连贯性。
模型在书籍数据上学到的
是如何维持长文本的\textbf{一致性和连贯性}：
这种能力在生成长篇回答、
撰写文章和创意写作中至关重要。

一个有趣的实验验证了这一点：
Touvron 等人在 LLaMA 1 的消融实验中发现，
将书籍数据的比例从 5\% 降到 0\%，
模型在短文本任务（如分类、NER）上几乎没有退化，
但在长文本生成任务（如故事续写、摘要生成）上
质量下降了 \textbf{8--12\%}：
证实了书籍数据对长文本能力的独特贡献。

\subsection{代码：推理能力的秘密武器}

\textbf{代码}是一个令人意外的关键数据源。
来自 GitHub、GitLab 等平台的开源代码不仅教会模型编程，
还能提升它的\textbf{通用推理能力}。
代码具有严格的逻辑结构——语法错误一个字符就会导致程序崩溃：
这种精确性迫使模型学习更严谨的推理模式。
Qwen3~\cite{qwen2025qwen3} 在 36T token 中
纳入了大量代码数据，正是为了利用这一效应。

代码数据的主要来源包括：

\begin{itemize}[noitemsep]
  \item \textbf{The Stack v1/v2}（BigCode / HuggingFace）：
        从 Software Heritage 存档（全球最大的源代码存档）
        中提取，覆盖 600+ 种编程语言。
        The Stack v2 包含约 67TB 的源代码。
        关键的设计决策是\textbf{许可证过滤}：
        只包含使用宽松开源许可证
        （如 MIT、Apache 2.0、BSD）发布的代码，
        排除了 GPL 等具有 copyleft 要求的许可证。
        % NEW REF: kocetkov2023stack = Kocetkov, Denis and Li, Raymond and others. The Stack: 3 TB of permissively licensed source code. TMLR, 2023.
  \item \textbf{GitHub 直接爬取}：
        许多团队直接从 GitHub API 获取代码。
        StarCoder 的数据流水线是一个优秀的参考：
        它对 GitHub 代码进行了多层过滤：
        去除自动生成的代码（如 \texttt{package-lock.json}）、
        去除过长或过短的文件、
        去除高比例的非 ASCII 字符、
        近似去重（代码的复制粘贴现象比文本更严重）。
        % NEW REF: li2023starcoder = Li, Raymond and others. StarCoder: may the source be with you! TMLR, 2023.
\end{itemize}

代码数据的处理还需要特别注意
\textbf{个人身份信息}（PII）的清理：
代码注释和配置文件中经常包含
开发者的邮箱、API 密钥、甚至密码。
这些必须在训练前被检测和移除。
StarCoder 的 PII 清理流水线使用了三层检测：
(1) 正则表达式匹配已知格式（邮箱、IP 地址、API 密钥模式）；
(2) NER 模型检测人名和组织名；
(3) 基于熵的检测器识别高随机性的字符串
（可能是密码或 token）。
检测到的 PII 被替换为占位符（如 \texttt{<EMAIL>}、\texttt{<KEY>}），
而非简单删除：
这保留了代码的结构完整性，
同时消除了隐私风险。

代码数据的\textbf{去重策略}也与文本数据不同。
代码的复制粘贴现象极为普遍：
GitHub 上大量的 fork 和 boilerplate 代码
导致代码数据的原始重复率高达 \textbf{60--70\%}
（远高于网页文本的 30--40\%）。
The Stack 使用了\textbf{精确去重}（文件级 SHA-256 哈希）
加\textbf{近似去重}（MinHash，阈值 0.7）的两层策略，
将代码数据量减少了约 50\%。
一个微妙的设计决策是：
对于 \texttt{LICENSE}、\texttt{README.md}
和 \texttt{.gitignore} 等“模板文件”，
采用\textbf{更激进}的去重（阈值 0.5），
因为这些文件在不同仓库间高度相似
但对模型学习编程能力几乎没有贡献。

\subsection{学术论文}

\textbf{学术论文}（来自 arXiv、PubMed、Semantic Scholar 等）
提供高质量的科学知识和数学推理内容。

\begin{itemize}[noitemsep]
  \item \textbf{arXiv}：包含约 200 万篇论文的 \LaTeX{} 源代码，
        是数学和科学推理数据的主要来源。
        直接使用 \LaTeX{} 源代码（而非 PDF 提取的文本）
        能保留公式的精确结构：
        这对模型学习数学推理至关重要。
  \item \textbf{S2ORC}（Semantic Scholar Open Research Corpus）：
        包含约 8100 万篇学术论文的全文和元数据，
        覆盖几乎所有学科。
  \item \textbf{PubMed / PMC}：生物医学领域的论文全文，
        约 3600 万篇摘要和 900 万篇全文。
\end{itemize}

PDF 到文本的转换是学术论文数据处理中的一大挑战。
PDF 是一种\textbf{面向排版}的格式，
不保留文档的逻辑结构。
表格、公式、多栏布局、脚注：
这些在 PDF 中的表示与人类阅读的理解完全不同。
让我们理解为什么这个问题如此困难：
PDF 将页面表示为\textbf{绘图指令}的序列：
“在坐标 (72, 648) 处用 12pt Times New Roman 绘制字符 `A'\thinspace”：
而不是“这是第 1 段的第一个字符”。
重建文档的逻辑结构
需要从这些低层绘图指令中
\textbf{逆向推断}段落边界、阅读顺序和语义结构。

2024--2025 年出现了几个重要的 PDF 解析工具：

\begin{itemize}[noitemsep]
  \item \textbf{Nougat}（Meta，2023）：
        使用视觉 Transformer（ViT + mBART decoder）
        将 PDF 页面的\textbf{图像}直接转换为 Markdown。
        这种“OCR + 理解”的端到端方法
        在数学公式上的准确率达到了 $\sim$90\%：
        远超传统的规则解析方法。
        但 Nougat 的速度较慢
        （每页约 2--3 秒），
        处理数百万篇论文需要大量 GPU 资源。
  \item \textbf{GROBID}：
        Java 实现的结构化论文解析器，
        使用 CRF（条件随机场）和规则引擎
        来识别标题、摘要、正文、引用等结构。
        速度快（每页 $<$0.1 秒），
        但在数学公式和复杂表格上表现较差。
  \item \textbf{Marker}（2024）：
        一个轻量级的开源工具，
        结合了 pdftotext 的文本提取
        和 LayoutLM 的布局理解，
        在速度和质量之间取得了较好的平衡。
  \item \textbf{MinerU}（上海 AI Lab，2024 年）：% NEW REF: mineru2024 = MinerU: An Open-Source Solution for Precise Document Content Extraction, 2024
        面向中文和多语言学术论文的开源解析工具，
        在中文论文（含中英混排、复杂表格）上的表现
        显著优于 Nougat 和 GROBID。
\end{itemize}

实践中的最佳策略是\textbf{优先使用源码}：
对 arXiv 上的论文直接使用 \LaTeX{} 源码
（通过 \texttt{tex2text} 或 \texttt{pandoc} 转换），
只对没有源码的论文才使用 PDF 解析。
arXiv 上约 80\% 的论文提供了 \LaTeX{} 源码：
这使得 arXiv 成为数学和科学知识
\textbf{质量最高}的数据来源，
远优于从 PDF 提取的版本。
对于 PubMed 等以 PDF 为主的来源，
Nougat + 后处理的组合是当前的最优方案，
但仍存在约 5--10\% 的公式解析错误率。

\subsection{对话与社交媒体}

\textbf{对话数据}（来自论坛、社交媒体等）
帮助模型学习自然的交流模式。
Reddit 是最常用的来源：
其 Pushshift 存档包含了数十年的帖子和评论，
线程结构的讨论数据对模型学习对话和论证能力特别有价值。
但 2023 年后 Reddit 限制了 API 访问，
使得新数据的获取变得困难。

Stack Overflow 和其他 Q\&A 平台也是重要来源：
问题和被采纳答案的配对
天然构成了高质量的问答训练数据。

\textbf{Twitter/X 数据}曾是对话数据的重要来源。
Twitter 的实时性和话题多样性使其成为训练模型
理解时事、俚语和非正式语言的关键语料。
但 2023 年 Elon Musk 收购 Twitter 并将其改名为 X 之后，
API 访问费用大幅提高（免费访问几乎被完全取消，
基础 API 月费从零上涨到 \$100--\$42000），
同时 Terms of Service 明确禁止将数据用于 AI 训练。
这使得合法获取 Twitter/X 数据用于模型训练
变得极其困难。
讽刺的是，xAI 的 Grok 模型被广泛认为使用了 X 平台数据：
但这种“平台所有者自用”的优势并不对外开放。
% NEW REF: twitter_api_2023 = Twitter API pricing changes and implications for AI research, 2023.

\textbf{Discord 和 Telegram} 等即时通讯平台
包含大量自然的多人对话数据：
这些对话具有快速轮换、非正式表达和多主题交织的特点，
与论坛式的结构化讨论（如 Reddit）形成互补。
但这些平台的数据获取面临\textbf{隐私}和\textbf{合规性}的双重挑战：
Discord 的 ToS 明确禁止对用户消息的批量爬取，
而 Telegram 的开放 API 虽然技术上可行，
但用户的隐私期望使得此类数据的使用
在欧盟 GDPR 等法律框架下存在法律风险。

\textbf{中文社交媒体}构成了中文 LLM 训练的重要数据来源。
微博提供了类似 Twitter 的短文本对话和舆论讨论数据，
但其 API 限制更为严格，且内容审查机制使得数据分布
与英文社交媒体存在系统性差异。
知乎（中国版 Quora）提供高质量的长篇问答数据，
其“回答”的平均长度和深度远超英文同类平台，
特别是在科技、人文和社会话题上。
小红书则以生活分享和消费决策为主，
提供了大量带有情感和经验性知识的文本：
这种类型的数据在传统的“百科全书+书籍+代码”训练数据组合中
往往缺失，但对模型的日常对话能力和情感理解至关重要。
百度贴吧、B 站评论区和豆瓣等平台
也各自覆盖了独特的话题和文化维度。

值得注意的是，2024--2025 年多个中文平台
开始明确限制 AI 训练的数据获取：
知乎在 2024 年修改了 robots.txt 禁止主流爬虫，
小红书在其用户协议中增加了关于 AI 训练的限制条款。
这些变化正在重塑中文 LLM 的数据供给格局，
也推动了中文合成数据生成的研究。
% NEW REF: zhihu_robots_2024 = Zhihu updates robots.txt to block AI crawlers, 2024.

\subsection{策划型数据集}

2023--2025 年，社区逐渐认识到
仅仅从单一来源提取数据是不够的：
需要\textbf{精心策划}的多源混合数据集。
代表性的工作包括：

\begin{itemize}[noitemsep]
  \item \textbf{RedPajama}（Together AI）：
        复现 LLaMA 的训练数据组成，
        包含 1.2T token，来自 Common Crawl、
        Wikipedia、书籍、arXiv、GitHub 和
        Stack Exchange。
  \item \textbf{SlimPajama}（Cerebras）：
        在 RedPajama 基础上做了更彻底的去重和清洗，
        将数据量从 1.2T 精简到 627B token：
        数据量减半但训练效果更好。
  \item \textbf{Dolma}（AI2）：
        Allen AI 为 OLMo 模型构建的 3T token 数据集，
        特点是\textbf{完全透明}：
        所有的数据处理步骤、过滤规则、
        混合比例都详细记录并公开发布。
        % NEW REF: soldaini2024dolma = Soldaini, Luca and others. Dolma: an Open Corpus of Three Trillion Tokens for Language Model Pretraining Research. arXiv preprint arXiv:2402.00159, 2024.
\end{itemize}

策划型数据集的演进揭示了一个重要趋势：
\textbf{数据透明度}正在成为开源 LLM 社区的核心价值观。
早期的数据集（如 C4、The Pile）
仅描述了数据来源和大致的清洗规则；
RedPajama 进一步公开了完整的数据组成和处理脚本；
而 Dolma 和 FineWeb 则达到了前所未有的透明度：
不仅公开了每一步处理的参数和随机种子，
还提供了中间产物和消融实验结果，
使得任何团队都可以精确复现其数据处理流程。
这种透明度对于\textbf{科学可复现性}至关重要：
如果我们不知道一个模型用了什么数据训练，
就无法真正理解它的能力和局限。

从 RedPajama 到 SlimPajama 的演进尤其说明了
\textbf{减量增质}的有效性。
SlimPajama 通过更激进的 MinHash 去重
（将阈值从 0.8 降到 0.7）
和更精细的质量过滤（增加了困惑度过滤和重复行检测），
将数据量从 1.2T 减少到 627B token：
但在相同训练 FLOP 下，模型性能反而提升了 3--5\%。
这再次印证了本章的核心论点：
在预训练数据工程中，质量的价值远大于数量。

\section{数据清洗流水线}

原始数据到训练数据之间，是一条复杂的清洗流水线。
让我们以 FineWeb~\cite{penedo2024fineweb} 的流水线为例来详细说明每个步骤。

\cref{fig:data-pipeline} 展示了一个典型的数据处理流水线的全景视图。
从原始的 Common Crawl 数据出发，
经过七个核心步骤后，
才能产出可以直接用于模型训练的 token 流。
每一步都会移除大量数据：
最终产出的高质量 token 通常只有原始数据量的 \textbf{10--15\%}。

\begin{figure}[htbp]
\centering
\begin{tikzpicture}[
  every node/.style={font=\small},
  stage/.style={
    rectangle, rounded corners=2pt, draw=#1!50, fill=#1!8,
    minimum width=24mm, minimum height=9mm,
    inner sep=5pt, align=center, line width=0.6pt,
    font=\small
  },
  arr/.style={-{Stealth[length=2.5pt,width=2.5pt]}, line width=0.6pt, figGray!60},
  vol/.style={font=\tiny, text=figGray!70, align=center},
  phase/.style={rounded corners=4pt, inner sep=6pt, line width=0pt},
]

% === Phase backgrounds (drawn first, behind everything) ===
\begin{scope}[on background layer]
  % Phase 1: Extraction & filtering (row 1, centered around nodes)
  \fill[figBlue!4, rounded corners=4pt]
    (-2.0, -0.85) rectangle (13.4, 1.35);
  % Phase 2: Curation (row 2, centered around nodes)
  \fill[figOrange!4, rounded corners=4pt]
    (-2.0, -4.05) rectangle (13.4, -1.85);
\end{scope}

% Phase labels
\node[font=\tiny\itshape, text=figBlue!50, anchor=north west] at (-1.8, 1.3) {Extraction \& Filtering};
\node[font=\tiny\itshape, text=figOrange!50, anchor=north west] at (-1.8, -1.9) {Curation \& Packaging};

% === Row 1: Extraction & Filtering ===
\node[stage=figBlue] (crawl) at (0, 0) {Common Crawl\\[-1pt]{\scriptsize (raw HTML)}};
\node[stage=figBlue] (extract) at (3.8, 0) {Text Extraction\\[-1pt]{\scriptsize trafilatura}};
\node[stage=figGreen] (lang) at (7.6, 0) {Lang ID\\[-1pt]{\scriptsize fastText}};
\node[stage=figOrange] (quality) at (11.4, 0) {Quality Filter\\[-1pt]{\scriptsize rules + classifier}};

% Row 1 arrows with volume annotations
\draw[arr] (crawl) -- (extract);
\node[vol, above=0.5mm of $(crawl.east)!0.5!(extract.west)$] {100s PB};
\draw[arr] (extract) -- (lang);
\draw[arr] (lang) -- (quality);

% === Row 2: Curation ===
\node[stage=figRed] (dedup) at (11.4, -3) {Deduplication\\[-1pt]{\scriptsize MinHash + LSH}};
\node[stage=figOrange] (domain) at (7.6, -3) {Domain Tagging\\[-1pt]{\scriptsize classification}};
\node[stage=figGreen] (mix) at (3.8, -3) {Data Mixing\\[-1pt]{\scriptsize domain ratios}};
\node[stage=figBlue] (tokenize) at (0, -3) {Tokenize \& Pack\\[-1pt]{\scriptsize BPE + sharding}};

% Connecting arrow between rows
\draw[arr] (quality.south) -- (dedup.north);
\node[vol, right=1mm of $(quality.south)!0.5!(dedup.north)$] {$-$60\%};

% Row 2 arrows with volume annotations
\draw[arr] (dedup) -- (domain);
\node[vol, above=0.5mm of $(dedup.west)!0.5!(domain.east)$] {$-$35\%};
\draw[arr] (domain) -- (mix);
\draw[arr] (mix) -- (tokenize);

% === Output ===
\node[stage=figGray, minimum width=28mm] (output) at (0, -5.2) {Training Shards\\[-1pt]{\scriptsize token sequences}};
\draw[arr] (tokenize.south) -- (output.north);
\node[vol, right=1mm of $(tokenize.south)!0.5!(output.north)$] {10--15\%};

\end{tikzpicture}
\caption{LLM 预训练数据流水线。原始 Common Crawl HTML 依次经过文本抽取、语言识别、
质量过滤（约去除 $60\%$）、多级去重（约去除 $35\%$）、领域分类、比例混合与
tokenization。最终高质量 token 流通常仅占原始数据量的 10--15\%。}
\label{fig:data-pipeline}
\end{figure}

\subsection{文本提取与规范化}

首先要从 HTML 中提取纯文本。
这不是简单的「去掉标签」：
需要保留有意义的结构（标题、段落、列表），
去除导航栏、页脚、侧边栏等「模板」文本。

\textbf{trafilatura} 是目前最流行的网页文本提取工具。
它使用基于规则和启发式的方法
来区分「正文内容」和「模板/导航」：
分析 DOM 树的结构、文本密度、
标签语义（\texttt{<article>} vs.\ \texttt{<nav>}）等。
在 FineWeb 的评估中，trafilatura 的提取质量
显著优于其他工具（如 jusText、readability）。

文本提取之后还需要做\textbf{规范化}：
统一 Unicode 编码（NFC/NFKC 规范化）、
移除控制字符和零宽空格、
统一引号和破折号的形式、
修复编码错误（mojibake）。
这些看似琐碎的步骤对后续的去重和质量过滤
有着不可忽视的影响：
如果两篇相同的文章使用了不同的 Unicode 编码形式，
去重算法可能无法检测到它们的重复。

\subsection{语言识别}

多语言清洗流水线的第一步是
判断每个文档的语言。
\textbf{fastText} 的语言分类器
是最广泛使用的工具：
它可以识别 176 种语言，
在句子级别的分类准确率超过 95\%。
fastText 的优势在于\textbf{极高的速度}：
它可以在单核 CPU 上以每秒数十万句的速度
进行语言分类，
这对于处理 PB 级的 Common Crawl 数据至关重要。

但 fastText 的\textbf{置信度分数}需要仔细校准。
FineWeb-2 的分析发现，
fastText 对高资源语言（英文、中文、法文）
的置信度通常在 0.95+ 范围，
但对低资源语言（如豪萨语、约鲁巴语）
的置信度中位数仅为 0.6--0.7：
因为训练 fastText 的语言标注数据
在低资源语言上本身就不足。
因此，FineWeb-2 为不同语言设定了
\textbf{差异化的置信度阈值}：
英文要求 $\geq 0.65$（因为数据量大，宁可严格）、
法文要求 $\geq 0.55$，
而豪萨语等低资源语言只要求 $\geq 0.35$
（否则会丢失过多本已稀缺的数据）。
这种\textbf{资源感知的阈值设计}
是多语言数据处理中一个容易被忽视但至关重要的细节。

语言识别的另一个微妙之处是
\textbf{混合语言}文档的处理。
许多网页（尤其是非英语国家的技术博客）
混合使用英语和本地语言。
简单地按文档级别分类会丢失这些有价值的混合语言内容。
FineWeb 采用的策略是：
对每个文档做段落级别的语言识别，
只要主要语言（$>$65\%）是目标语言就保留。

\textbf{CLD3}（Compact Language Detector v3，Google 开源）
是 fastText 的主要替代方案。
CLD3 使用了一个小型神经网络
（而非 fastText 的线性分类器），
在混合语言文本上的表现优于 fastText：
但速度约慢 3--5 倍。
DCLM 采用了\textbf{双分类器投票}策略：
只有 fastText \textbf{和} CLD3 都同意某文档属于目标语言时才保留。
这种保守策略将语言误判率从约 5\% 降低到 $<$1\%，
代价是约 8\% 的召回率损失：
对于数据量充裕的英文场景，这是一个合理的取舍；
但对于低资源语言，单分类器加低阈值更为合适。

\textbf{脚本检测}（script detection）是语言识别的一个补充步骤，
在 FineWeb-2 的多语言处理中扮演重要角色。
某些语言使用独特的书写系统
（如天城文用于印地语、格鲁吉亚文字、缅甸文等），
可以通过 Unicode 字符范围检测来快速确定：
这比运行完整的语言分类器快 100 倍以上。
FineWeb-2 使用“先脚本、后语言”的\textbf{两步策略}：
先通过 Unicode 脚本检测快速分桶
（将天城文脚本的文档放入“印度语系”桶），
然后只在桶内运行 fastText 来区分具体语言
（印地语 vs 马拉地语 vs 尼泊尔语）。
这种层级策略将语言识别的总计算量减少了约 40\%。

\subsection{去重}

互联网上的内容重复非常严重：
同一篇新闻可能被数百个网站转载。
如果不去重，模型会在重复内容上浪费大量训练计算，
更严重的是，它可能直接记住这些文本而非学习泛化的模式。

研究表明，Common Crawl 中约 \textbf{30--40\%}
的内容是近似重复的。
Lee 等人（2022）的研究发现，
去重可以带来与增加 3--4 倍训练数据
等价的模型质量提升：
这使得去重成为数据处理中
投入产出比最高的步骤之一。
% NEW REF: lee2022dedup = Lee, Katherine and others. Deduplicating Training Data Makes Language Models Better. In ACL, 2022.

\subsubsection{精确去重}

\textbf{精确去重}是最简单的形式：
对每个文档计算哈希值（如 SHA-256），
哈希值相同的文档被视为完全重复。

\textbf{URL 去重}是精确去重的一个变体：
对于来自同一 URL 的多次爬取
（Common Crawl 的不同月度快照可能爬取同一网页），
只保留最新的版本。
FineWeb 在处理 96 个快照时，
URL 去重就移除了约 40\% 的数据。

精确去重的局限很明显：
它无法检测到只有微小差异的文档：
例如，同一篇文章被不同网站转载时
可能添加了不同的广告文本或页脚。

\subsubsection{模糊去重：MinHash + LSH}

\textbf{模糊去重}（fuzzy deduplication）使用
MinHash + LSH（局部敏感哈希）
找到高度相似但不完全相同的文档。
这是整个数据处理流水线中最精巧的算法之一，
值得详细理解。

\textbf{第一步：Shingling（构建 $n$-gram 集合）。}
将每个文档表示为其所有连续 $n$ 个词组成的片段
（称为 shingle）的集合。
例如，对于文档 “the cat sat on the mat”，
使用 3-gram shingling：

\begin{center}
\small
$\{$“the cat sat”, “cat sat on”, “sat on the”, “on the mat”$\}$
\end{center}

FineWeb 使用 5-gram shingling。
$n$ 的选择是一个权衡：
$n$ 太小会导致过多的误报（不相关的文档也可能共享短片段），
$n$ 太大会导致过多的漏报（相似文档可能没有足够长的共享片段）。

\textbf{第二步：Jaccard 相似度。}
两个文档的相似程度用 Jaccard 系数衡量：
\begin{equation}
  J(A, B) = \frac{|A \cap B|}{|A \cup B|}
\end{equation}
其中 $A$ 和 $B$ 是两个文档的 shingle 集合。
如果 $J(A, B) > 0.8$，我们认为这两个文档是近似重复的。

但直接计算所有文档对的 Jaccard 相似度需要 $O(n^2)$ 次比较：
对于数十亿个文档，这是不可行的。

\textbf{第三步：MinHash 签名。}
MinHash 用一组 $k$ 个随机哈希函数
将每个文档映射为一个长度为 $k$ 的「签名」。
具体做法是：对每个哈希函数 $h_i$，
计算文档所有 shingle 在 $h_i$ 下的哈希值，
取\textbf{最小值}作为签名的第 $i$ 个分量：
\begin{equation}
  \text{sig}_i(A) = \min_{s \in A} h_i(s)
\end{equation}

MinHash 的关键数学性质是：
\begin{equation}
  \Pr[\text{sig}_i(A) = \text{sig}_i(B)] = J(A, B)
\end{equation}
即两个签名在第 $i$ 个位置相同的概率
\textbf{恰好等于}它们 shingle 集合的 Jaccard 相似度。

\begin{whybox}[为什么 MinHash 能近似 Jaccard？]
这个等式的证明非常优雅。
考虑集合 $A \cup B$ 中的所有 shingle。
每个 shingle 属于三种情况之一：
(a) 仅在 $A$ 中，(b) 仅在 $B$ 中，(c) 在 $A \cap B$ 中。
一个随机哈希函数 $h$ 给每个 shingle 一个随机值，
$\min$ 操作选中的那个 shingle
等可能是 $A \cup B$ 中的任何一个。
如果它恰好落在 $A \cap B$ 中，
则 $\text{sig}(A) = \text{sig}(B)$。
这发生的概率恰好是 $|A \cap B| / |A \cup B| = J(A, B)$。

使用 $k$ 个独立的哈希函数，
签名匹配的比例就是 $J(A, B)$ 的一个无偏估计，
误差随 $k$ 增大而减小（标准差为 $O(1/\sqrt{k})$）。
FineWeb 使用 $k = 128$，
意味着 Jaccard 相似度的估计误差约为 $\pm 0.09$。
\end{whybox}

因此，我们只需比较签名（$k$ 个整数），
不需要比较原始的 shingle 集合（可能有数万个元素）。
但即使比较签名，$O(n^2)$ 的文档对数仍然太多。

\textbf{第四步：LSH（Locality-Sensitive Hashing）。}
LSH 将长度为 $k$ 的签名分成 $b$ 个 band，
每个 band 包含 $r$ 个值（$k = b \times r$）。
对每个 band，对其 $r$ 个值再做一次哈希，
只有在\textbf{至少一个 band 上完全匹配}的文档对
才会被标记为候选相似对，然后做精确比较。

通过调整 $b$（band 数量）和 $r$（每个 band 的行数），
可以精确控制「召回率 vs.\ 误报率」的权衡。
两个 Jaccard 相似度为 $s$ 的文档被标记为候选对的概率是：
\begin{equation}
  P(\text{candidate} \mid J = s) = 1 - (1 - s^r)^b
\end{equation}

这个函数呈 S 形曲线。
当 $b = 20, r = 5$（即 $k = 100$）时：
\begin{itemize}[noitemsep]
  \item $s = 0.8$ 的文档对被检测到的概率 $\approx 0.997$（几乎全部召回）
  \item $s = 0.5$ 的文档对被检测到的概率 $\approx 0.47$（约一半误报）
  \item $s = 0.2$ 的文档对被检测到的概率 $\approx 0.006$（几乎不误报）
\end{itemize}

FineWeb 使用 $b = 14, r = 8$（$k = 112$），
相似度阈值约 0.75：
在去重召回率和计算成本之间取得了较好的平衡。

整个 MinHash + LSH 流水线将 $O(n^2)$ 的所有文档对比较
降低到近似线性的复杂度：
这使得在数十亿文档上做模糊去重变得实际可行。

\subsubsection{去重的定量效果}

去重在不同数据集上的效果差异巨大，
了解这些数字有助于设定合理的预期。

\begin{center}\small
\renewcommand{\arraystretch}{1.3}
\begin{tabular}{lrrr}
\toprule
\rowcolor{figLightGray}
\textbf{数据集} & \textbf{去重前} & \textbf{去重后} & \textbf{移除比例} \\
\midrule
FineWeb（URL 去重） & 96 快照原始数据 & -- & $\sim$40\% \\
FineWeb（MinHash 去重） & URL 去重后 & 15T token & $\sim$25\% \\
SlimPajama（全流程） & 1.21T token & 627B token & 48\% \\
RedPajama v2（MinHash） & -- & -- & $\sim$30\% \\
RefinedWeb（全流程） & -- & 5T token & $\sim$35\% \\
\bottomrule
\end{tabular}
\end{center}

从 RedPajama 到 SlimPajama 的演进
提供了一个精确的\textbf{去重激进度对比}。
RedPajama 使用 MinHash 相似度阈值 0.8
（Jaccard $> 0.8$ 的文档对被视为重复），
而 SlimPajama 将阈值降低到 0.7：
这意味着即使两个文档只有 70\% 的内容相同，
也会被标记为近似重复。
降低阈值的代价是\textbf{误删率上升}：
一些内容相关但并不真正重复的文档
（如同一事件的不同报道）也会被移除。
但 Cerebras 团队的消融实验表明，
在相同训练 FLOP 下，
更激进的去重带来了 3--5\% 的下游任务提升，
证明了“宁可多删、不可遗漏”的原则
在数据充足的场景下是正确的。

MinHash 的\textbf{参数选择}是一门精细的工程。
核心参数有三个：
shingle 大小 $n$、哈希函数数量 $k$、
以及 LSH 的 band/row 配置 $(b, r)$。
它们之间的关系是 $k = b \times r$，
且阈值近似为 $(1/b)^{1/r}$。
以下是几个主要项目的参数选择对比：

\begin{center}\small
\renewcommand{\arraystretch}{1.3}
\begin{tabular}{lccccc}
\toprule
\rowcolor{figLightGray}
\textbf{项目} & $n$ (shingle) & $k$ (hashes) & $b$ (bands) & $r$ (rows) & 阈值 \\
\midrule
FineWeb & 5 & 112 & 14 & 8 & $\sim$0.75 \\
SlimPajama & 5 & 128 & 9 & 13 & $\sim$0.70 \\
RefinedWeb & 5 & 128 & 20 & 5 & $\sim$0.80 \\
Dolma & 5 & 128 & 14 & 8 & $\sim$0.75 \\
\bottomrule
\end{tabular}
\end{center}

注意 SlimPajama 的 $(b=9, r=13)$ 配置：
更少的 band、更多的 row per band：
这使得 S 曲线的“拐点”左移到约 0.70，
从而在更低的相似度下就开始积极检测重复。
代价是更高的计算量（每个 band 中更多的行需要全部匹配），
但 Cerebras 认为这个代价对数据质量的改善是值得的。

去重操作的\textbf{计算规模}不容低估。
FineWeb 处理 96 个 Common Crawl 快照的 MinHash 去重
需要在分布式集群上运行，
每个快照约包含 30 亿个文档。
MinHash 签名的存储量为 $30 \times 10^9 \times 112 \times 4$ bytes
$\approx 13$\,TB（假设每个哈希值用 32 位整数存储）。
LSH 索引的构建和查询需要额外的临时存储。
HuggingFace 使用了约 \textbf{50 万 CPU 小时}
仅用于 FineWeb 的去重步骤：
这几乎是整个数据处理流水线中
计算成本最高的单一步骤。

\subsubsection{后缀数组去重}

除了 MinHash + LSH 这种文档级别的去重，
还有一种更精细的去重方式：
\textbf{后缀数组去重}（suffix array deduplication）。
Lee 等人（2022）提出了这种方法，
LLaMA 团队~\cite{touvron2023llama}也采用了它。

后缀数组去重在\textbf{子字符串}级别工作：
它找到语料中所有重复出现的长子字符串
（通常 $\geq$ 50 个字符），
然后只保留一个副本，其余删除。
这种方法可以检测到
\textbf{段落级别的重复}：
例如，许多网页可能共享相同的免责声明、
版权声明或 cookie 通知。

后缀数组的构建可以在 $O(n)$ 时间内完成
（使用 SA-IS 等算法），
但需要的内存约为文本大小的 5--8 倍。
对于 TB 级别的语料，
这需要分布式实现和外部排序。

\subsubsection{去重的层次}

总结来说，一个完整的去重流水线通常包含多个层次：

\begin{center}\small
\renewcommand{\arraystretch}{1.3}
\begin{tabular}{lll}
\toprule
\rowcolor{figLightGray}
\textbf{层次} & \textbf{方法} & \textbf{检测目标} \\
\midrule
URL 去重       & 精确匹配         & 同一 URL 的多次爬取 \\
文档精确去重   & SHA-256 哈希     & 完全相同的文档 \\
文档模糊去重   & MinHash + LSH    & 高度相似的文档 \\
段落/子串去重  & 后缀数组         & 共享的段落或模板 \\
\bottomrule
\end{tabular}
\end{center}

\subsection{质量过滤}

去重之后，仍然有大量低质量文本需要清除。
质量过滤是整个流水线中最需要「工程直觉」的环节：
什么是「高质量」并没有一个公认的数学定义。

\subsubsection{基于规则的启发式过滤}

最直接的方法是设计一组规则来过滤明显的低质量内容。
以 Gopher 模型（DeepMind）的过滤规则为代表：

\begin{itemize}[noitemsep]
  \item \textbf{文档长度}：过短（$<$ 50 词）或过长（$>$ 100K 词）的文档被移除。
  \item \textbf{平均词长}：平均词长 $<$ 3 或 $>$ 10 的文档被移除
        （可能是乱码或非自然语言文本）。
  \item \textbf{特殊字符比例}：如果文档中
        非字母数字字符超过 30\%，被移除。
  \item \textbf{重复行比例}：如果超过 30\% 的行是重复的，被移除
        （可能是导航菜单或表格的残留）。
  \item \textbf{大写字母比例}：如果超过 70\% 的字母是大写的，被移除
        （可能是全大写的标题或垃圾邮件）。
  \item \textbf{停用词比例}：如果英语停用词（“the”、“a”、“is”……）
        的比例低于 5\%，被移除
        （可能是关键词列表而非自然文本）。
  \item \textbf{特征字符串检测}：包含 “lorem ipsum”、
        “javascript is disabled”、“cookie policy”
        等特征字符串的文档被移除。
\end{itemize}

这些规则看似简单粗暴，但效果出奇地好：
它们可以在几乎不损失高质量内容的情况下
移除大量明显的垃圾。

值得深入了解的是这些规则背后的\textbf{定量阈值}设计：
它们并非拍脑袋决定的，
而是经过大量消融实验确定的最优工作点。
以 Gopher 的规则为基础，
FineWeb~\cite{penedo2024fineweb} 和 DCLM 团队
进一步精细化了这些阈值：

\begin{center}\small
\renewcommand{\arraystretch}{1.3}
\begin{tabular}{lcc}
\toprule
\rowcolor{figLightGray}
\textbf{过滤规则} & \textbf{Gopher 阈值} & \textbf{FineWeb 优化阈值} \\
\midrule
最小文档长度 & 50 词 & 65 词 \\
最大文档长度 & 100,000 词 & 100,000 词 \\
平均词长范围 & [3, 10] & [3, 10] \\
最大特殊字符比例 & 30\% & 20\% \\
最大重复行比例 & 30\% & 20\% \\
最小停用词比例 & 5\% & 6.5\% \\
最大大写字母比例 & 70\% & 60\% \\
重复 2-gram 比例上限 & -- & 20\% \\
重复 3-gram 比例上限 & -- & 15\% \\
重复 4-gram 比例上限 & -- & 10\% \\
\bottomrule
\end{tabular}
\end{center}

FineWeb 引入的\textbf{重复 $n$-gram 检测}是一个重要的创新。
与“重复行”检测不同，$n$-gram 重复检测能够捕捉到
更隐蔽的低质量模式，例如 SEO 优化的网页
经常在正文中反复插入相同的关键词短语，
每次出现的位置不同（因此不是“重复行”），
但 $n$-gram 频率异常高。
具体实现是：对文档中所有 $n$-gram 计数，
如果出现超过一次的 $n$-gram 覆盖的 token 数
占总 token 数的比例超过阈值，则判定为低质量。
FineWeb 在实验中发现，
\textbf{2-gram 重复检测对移除 SEO 垃圾最有效}，
而 \textbf{4-gram 重复检测对移除模板化内容最有效}
（如法律声明、Cookie 通知的多次出现）。

阈值的选择对最终数据集质量有\textbf{非线性}的影响。
DCLM 的消融实验揭示了一个“甜蜜区”现象：
以停用词比例为例，
将阈值从 5\% 提高到 6.5\%
在下游基准（HellaSwag、ARC）上带来了 0.8\% 的提升；
但进一步提高到 8\%
反而导致 0.3\% 的退化：
因为开始误删正常的技术文档
（代码教程中停用词比例天然偏低）。
这种非单调性使得阈值调优不能简单地“越严格越好”，
而必须在每个维度上找到最优工作点。

\textbf{多语言阈值的适配}是另一个非平凡的问题。
上述阈值大多针对英文设计：
它们对中文、日文等非空格分隔的语言不直接适用。
例如，“平均词长”的概念在中文中没有意义
（中文词之间没有空格），
“停用词比例”需要使用语言特定的停用词表。
FineWeb-2 为每种语言维护了一组\textbf{独立的阈值配置}，
通过在小规模标注数据上的交叉验证来确定：
这意味着 1000+ 种语言需要 1000+ 套不同的过滤参数。
对于标注数据极少的低资源语言，
FineWeb-2 采用了\textbf{语系聚类}策略：
将同一语系的语言（如班图语系的多种非洲语言）
共享一组阈值，并在此基础上做微调。
% NEW REF: penedo2025fineweb2_thresh = FineWeb-2 uses per-language quality thresholds derived from cross-validation on annotated samples.

\subsubsection{基于模型的质量过滤}

更精细的过滤使用\textbf{机器学习模型}来判断文档质量。

\textbf{困惑度过滤}（Perplexity filtering）：
使用一个在高质量文本（如 Wikipedia）上
训练的小型语言模型（如 KenLM）
计算每个文档的困惑度。
直觉上，高质量的文本（如维基百科文章）
在语言模型眼中是「正常」的，困惑度较低；
而低质量的文本（如乱码或机器生成的垃圾）
具有异常高的困惑度。
CCNet（Meta 的 Common Crawl 清洗工具）
就是基于这种方法，
将文档按困惑度分为「头部」（高质量）、
「中部」和「尾部」（低质量）三个等级。
% NEW REF: wenzek2020ccnet = Wenzek, Guillaume and Lachaux, Marie-Anne and Conneau, Alexis and Chaudhary, Vishrav and Guzman, Francisco and Joulin, Armand and Grave, Edouard. CCNet: Extracting High Quality Monolingual Datasets from Web Crawl Data. In LREC, 2020.

困惑度过滤的\textbf{定量阈值}选择至关重要。
CCNet 将文档按困惑度分为三个等级：
“头部”（困惑度 $< 230$，约占文档总数的 33\%）、
“中部”（$230 \leq$ 困惑度 $< 560$，约 33\%）、
“尾部”（困惑度 $\geq 560$，约 33\%）。
这些阈值基于 5-gram KenLM 模型在 Wikipedia 上训练后的分布。
LLaMA 选择了仅使用“头部”文档：
这个简单的阈值选择在 2023 年被证明
足以构建有竞争力的训练数据。

但这个三等分阈值并非最优。
FineWeb 的消融实验表明，
将困惑度阈值设为 \textbf{345}（而非 CCNet 的 230）
在下游任务上表现更好：
因为 230 的阈值过于严格，
排除了大量质量尚可但写作风格与 Wikipedia 不同的文档。
DCLM 进一步引入了\textbf{双侧阈值}：
不仅过滤高困惑度文档（$> 400$，可能是乱码），
也过滤\textbf{极低}困惑度文档（$< 10$，
可能是简单的重复文本或自动生成的模板内容）。
这种双侧过滤的思路颇具洞察力：
困惑度过低说明文本缺乏信息量，
模型从中几乎学不到新东西。

困惑度过滤还有一个更深层的陷阱：
它倾向于保留“类似 Wikipedia”的文本，
可能过度偏向百科全书式的写作风格，
而丢弃同样有价值的其他风格
（如技术教程、个人博客、论坛讨论）。
这种\textbf{风格偏见}会导致训练出的模型
在回答问题时过于正式和百科化，
缺乏自然对话的灵活性。
一种缓解策略是为不同\textbf{领域}
训练不同的困惑度模型：
例如，在 Stack Overflow 上训练的 KenLM
用于过滤技术文档，
在 Reddit 上训练的 KenLM
用于过滤对话数据。
这种“领域自适应困惑度过滤”
在 FineWeb-2 的多语言处理中被广泛采用。

\textbf{分类器过滤}：
GPT-3~\cite{brown2020gpt3} 的训练
使用了一个二分类器来判断文档是否「高质量」：
正样本来自 Wikipedia 和 WebText（Reddit 上高赞帖子链接的网页），
负样本来自原始 Common Crawl。
LLaMA~\cite{touvron2023llama} 也采用了类似的策略。

FineWeb~\cite{penedo2024fineweb} 的关键创新之一
是使用了多轮、多方法的质量过滤，
包括基于规则的启发式过滤
和基于模型的打分过滤。
特别是 FineWeb-Edu 子集使用了一个
专门训练的\textbf{教育质量分类器}：
这个分类器评估文档是否具有教育价值，
而不仅仅是「语言质量是否高」。
结果表明，1.3T token 的 FineWeb-Edu
在多项基准测试上的训练效果
甚至优于使用全量 15T token 的 FineWeb。
这是「质量 $>$ 数量」的最有力的实证之一。

\subsubsection{安全过滤}

\textbf{NSFW/毒性过滤}是数据清洗的重要一环。
模型如果在色情、暴力、仇恨言论等内容上训练，
可能会在推理时生成类似的有害内容。

常用的工具包括：
\begin{itemize}[noitemsep]
  \item \textbf{Jigsaw Perspective API}：
        Google 开发的毒性评分工具，
        可以为文本标注「毒性」、「侮辱性」、
        「威胁性」等维度的分数。
  \item \textbf{自训练分类器}：
        使用少量标注数据微调一个小模型（如 fastText）
        来检测 NSFW 和毒性内容。
        FineWeb 使用了这种方法，
        对检测到的有害内容进行移除或降权。
\end{itemize}

\textbf{PII 移除}（个人身份信息）也是安全过滤的一部分：
\begin{itemize}[noitemsep]
  \item 基于正则表达式的检测：
        邮箱地址、电话号码、身份证号、信用卡号等
        具有固定格式，可以用正则表达式匹配和移除。
  \item 基于命名实体识别（NER）的检测：
        人名、地址等没有固定格式的 PII
        需要使用 NER 模型来识别。
\end{itemize}

\section{数据配比：一门精细的艺术}

确定了数据来源和清洗流水线之后，
还有一个关键问题：不同类型的数据应该以什么比例混合？

\subsection{经验法则与实验结果}

这并不是一个简单的优化问题。
LLaMA 3~\cite{dubey2024llama3} 的训练数据配比经过了大量实验。
据社区根据技术报告推测的数据配比大致为：
约 50\% 网页文本、25\% 代码、10\% 学术论文、
10\% 书籍、5\% 数学数据（官方未公布精确比例）。

一些令人惊讶的发现：
\begin{itemize}[noitemsep]
  \item 增加代码数据的比例可以提升模型在\textbf{非代码任务}
        （如数学推理和逻辑推理）上的表现。
  \item 增加数学数据的比例在超过某个阈值后
        反而会降低模型的通用能力。
  \item 高质量数据可以被重复使用若干次（通常 2--4 epochs）
        而不会导致严重的过拟合。
\end{itemize}

代码数据提升推理能力这一发现值得深入解释。
代码具有几个独特的特征使其成为推理能力的催化剂。
首先，代码有\textbf{严格的逻辑结构}：
一个括号没有匹配就会导致语法错误，
一个逻辑条件写错就会导致程序行为完全不同。
这种精确性迫使模型学习更严谨的逻辑推理。
其次，代码包含大量\textbf{抽象和组合}：
函数调用、类继承、模块导入：
这些正是推理能力的核心组成部分。
第三，代码的\textbf{执行结果是确定的}：
不像自然语言那样模糊，
代码的正确性可以通过执行来验证，
为模型提供了更清晰的学习信号。

\subsection{真实模型的配比参考}

不同模型的数据配比策略差异很大，
反映了各团队的设计哲学和优势领域：

\begin{center}\small
\renewcommand{\arraystretch}{1.3}
\begin{tabular}{lrccccc}
\toprule
\rowcolor{figLightGray}
\textbf{模型} & \textbf{总量} & \textbf{网页} & \textbf{代码} & \textbf{学术/书籍} & \textbf{数学} & \textbf{其他} \\
\midrule
LLaMA 3    & 15T  & 50\% & 25\% & 15\% & 5\%  & 5\%  \\
DeepSeek-V3 & 14.8T & 40\% & 30\% & 10\% & 10\% & 10\% \\
Qwen3      & 36T  & --   & 大量 & --   & 大量 & --   \\
LLaMA 4    & 40T  & --   & 大量 & --   & --   & 多模态 \\
\bottomrule
\end{tabular}
\end{center}

（各模型的详细配比大多未完全公开，
上表中的百分比为根据技术报告的描述推断。
LLaMA 4 在 40T token 中包含了 200 种语言的数据，
且原生支持图像输入，
暗示多模态数据在训练配比中占据了可观的份额。）

值得注意的是，DeepSeek-V3~\cite{deepseekv3}
将代码占比提高到约 30\%：
这可能是其在数学和推理任务上表现优异的原因之一。
而 LLaMA 4 的 40T token 训练数据中覆盖了 200 种语言：
这是迄今为止语言覆盖最广的 frontier 模型，
反映了 Meta 的全球化部署策略。

\subsection{经典数据集的领域组成}

在讨论自动配比优化之前，
先回顾几个经典数据集的领域组成，
了解“前人的智慧”是理解配比优化的出发点。

\textbf{The Pile}（EleutherAI，2021 年）
是最早的精心策划型数据集之一，
其 800B token 的组成如下：

\begin{center}\small
\renewcommand{\arraystretch}{1.3}
\begin{tabular}{lr}
\toprule
\rowcolor{figLightGray}
\textbf{数据源} & \textbf{占比} \\
\midrule
Pile-CC（Common Crawl 子集） & 52.5\% \\
PubMed Central & 14.4\% \\
Books3 & 12.1\% \\
OpenWebText2 & 8.4\% \\
ArXiv & 4.6\% \\
GitHub & 3.3\% \\
Wikipedia & 2.0\% \\
其他（Stack Exchange、FreeLaw 等） & 2.7\% \\
\bottomrule
\end{tabular}
\end{center}

The Pile 的设计哲学是\textbf{最大化多样性}：
它刻意纳入了法律文本（FreeLaw）、
专利文本（USPTO Backgrounds）、
数学证明（DM Mathematics）等冷门领域，
即使这些领域的数据量很小。
这种“宁缺毋偏”的设计思想
在后续被证明对模型的\textbf{长尾知识覆盖}
有重要贡献：
The Pile 上训练的 GPT-Neo 系列模型
在法律和科学推理任务上
表现优于数据量更大但领域更单一的竞品。
% NEW REF: gao2021pile = Gao, Leo and others. The Pile: An 800GB Dataset of Diverse Text for Language Modeling. arXiv preprint arXiv:2101.00027, 2021.

\textbf{SlimPajama} 的领域组成则反映了“精简”策略：

\begin{center}\small
\renewcommand{\arraystretch}{1.3}
\begin{tabular}{lrr}
\toprule
\rowcolor{figLightGray}
\textbf{数据源} & \textbf{RedPajama} & \textbf{SlimPajama} \\
\midrule
Common Crawl & 878B (73.2\%) & 521B (83.1\%) \\
C4 & 175B (14.6\%) & -- (合并到 CC) \\
GitHub & 59B (4.9\%) & 35B (5.6\%) \\
Wikipedia & 24B (2.0\%) & 17B (2.7\%) \\
Books & 26B (2.2\%) & 26B (4.1\%) \\
ArXiv & 28B (2.3\%) & 24B (3.8\%) \\
Stack Exchange & 10B (0.8\%) & 4B (0.6\%) \\
\midrule
\textbf{总计} & \textbf{1200B} & \textbf{627B} \\
\bottomrule
\end{tabular}
\end{center}

从 RedPajama 到 SlimPajama，
Common Crawl 的数据量减少了 40\%（主要因为更激进的去重），
但其在整体中的\textbf{占比反而上升}：
从 73\% 到 83\%。
这是因为 Common Crawl 中重复内容的比例最高，
去重移除的主要是 CC 数据；
而 GitHub、ArXiv 等相对干净的数据源
去重后保留率更高。
值得注意的是，SlimPajama 中\textbf{书籍数据的保留率最高}
（26B $\to$ 26B，几乎没有被去重移除），
这符合直觉——书籍文本经过专业编辑，重复率本就极低。

\subsection{自动数据配比：DoReMi 与后继者}

手动调整数据配比需要大量的试错实验，
成本极高。
2023 年，Xie 等人提出了 \textbf{DoReMi} 算法，
试图用自动化的方式找到最优配比。
% NEW REF: xie2023doremi = Xie, Sang Michael and Santurkar, Shibani and Ma, Tengyu and Liang, Percy. DoReMi: Optimizing Data Mixtures Speeds Up Language Model Pretraining. In NeurIPS, 2023.

DoReMi 的核心思想精妙而实用。
它分两步走：
\textbf{第一步}，在均匀配比上训练一个小型“参考模型”
（reference model，通常 280M 参数、几十 B token）；
\textbf{第二步}，训练一个同规模的“代理模型”（proxy model），
但训练过程中动态调整各领域的采样权重：
对“代理模型比参考模型差得最多”的领域加大采样。
这本质上是\textbf{分布式鲁棒优化}
（Distributionally Robust Optimization, DRO）
的在线版本：
它不是优化“平均性能最好”，
而是优化“最差领域也不太差”。

具体地，令 $\ell_d^{\text{ref}}$ 和 $\ell_d^{\text{proxy}}$
分别是参考模型和代理模型在领域 $d$ 上的损失。
DoReMi 使用\textbf{指数梯度}算法
在训练过程中动态更新每个领域的权重：
\begin{equation}
  w_d^{(t+1)} \propto w_d^{(t)} \cdot
  \exp\!\bigl(\eta \cdot [\ell_d^{\text{proxy},(t)} - \ell_d^{\text{ref}}]\bigr)
\end{equation}
当代理模型在某个领域的损失远高于参考模型时
（即该领域“学得最差”），
该领域的权重指数级上升，
更多的训练 batch 被分配给它。
这种\textbf{自适应上采样}策略
确保了所有领域都获得足够的训练信号。

实验表明，DoReMi 找到的配比
在下游任务上的表现
比均匀配比提升了约 6\%，
而且这种提升在不同规模的模型上都成立。
一个出人意料的发现是：
DoReMi 将 \textbf{Wikipedia 的权重大幅降低}
（从均匀配比的 $\sim$15\% 降到 $\sim$3\%），
因为 Wikipedia 的文本质量很高、分布简单，
模型很快就能学好，
继续在上面训练的边际收益很低；
相反，它将\textbf{代码和数学数据的权重提高}了 2--3 倍：
因为这些领域结构复杂，模型需要更多的训练才能掌握。
这个发现挑战了“高质量数据就应该多用”的朴素直觉：
正确的策略是“在模型最需要的领域多用”。

\subsubsection{RegMix：将配比优化建模为回归问题}

2025 年（ICLR），Liu 等人提出了 \textbf{RegMix}，
将数据配比问题重新建模为一个\textbf{回归任务}。
% NEW REF: liu2025regmix = Liu, Qian and Zheng, Xiaosen and Muennighoff, Niklas and Zeng, Guangtao and Dou, Longxu and Pang, Tianyu and Jiang, Jing and Lin, Min. RegMix: Data Mixture as Regression for Language Model Pre-training. In ICLR, 2025.

RegMix 的方法分三步：
(1) 在不同的数据配比上训练大量小模型（512 个 1M 参数的模型，各用 1B token）；
(2) 将每个配比作为输入、训练后的模型性能作为输出，
拟合一个回归模型来预测未尝试配比的性能；
(3) 用回归模型预测的最优配比来训练大规模模型。

RegMix 的关键优势在于\textbf{效率}：
它仅使用 DoReMi 约 10\% 的计算资源，
就能在 100B token 训练的 7B 模型上
匹配或超越 DoReMi 的表现。
一个令人意外的发现是：
回归分析表明，\textbf{网页语料}
（而非通常被认为更高质量的 Wikipedia）
与下游性能有最强的正相关：
这挑战了「高质量数据 = 类 Wikipedia 数据」的朴素假设。

\subsubsection{MDE：混合数据专家模型}

Google DeepMind 在 2025 年提出了
\textbf{Mixture of Data Experts}（MDE）方法，
进一步提升了配比优化的精度。
MDE 为每种数据领域训练一个小型「专家模型」，
然后通过这些专家模型的加权组合
来高效近似任意配比下的模型损失。
这种近似作为回归模型的额外特征，
使得配比预测的准确性显著提升。
% NEW REF: belenki2025mde = Belenki, Lior and Agarwal, Alekh and Shi, Tianze and Toutanova, Kristina. Optimizing Pre-Training Data Mixtures with Mixtures of Data Expert Models. arXiv preprint arXiv:2502.15950, 2025.

MDE 相对于前辈方法（DoReMi、RegMix）的核心改进
在于它构建了更精确的\textbf{损失预测模型}。
DoReMi 使用一个小型代理模型来估计不同配比的效果，
但代理模型本身的训练只能探索有限的配比空间：
它在本质上是一个单点估计，对未尝试过的配比区域泛化能力有限。
RegMix 通过训练 512 个微型模型来构建回归模型，
但这些模型的规模（1M 参数、1B token）
与目标模型（7B+ 参数、100B+ token）之间的差距
可能导致 scaling 行为的系统性偏差。
MDE 的“专家组合”策略则提供了一种理论上有据的近似：
基于混合模型损失可以被领域专家损失的加权组合
近似这一关键假设，
MDE 可以在不实际训练大量混合配比模型的情况下
“合成”任意配比的损失估计。

在计算成本方面，MDE 训练 $K$ 个领域专家模型的开销
约为 DoReMi 的 $K/N$ 倍
（$K$ 为数据领域数量，$N$ 为 DoReMi 的代理模型大小比例）。
对于典型的 5--10 个数据领域设置，
MDE 的总计算量约为 DoReMi 的 30--50\%，
同时在配比预测的 Kendall's $\tau$ 排序相关性上
达到了 0.85+（DoReMi 约 0.6、RegMix 约 0.75）。
这意味着 MDE 推荐的配比方案
在实际大模型训练中更可靠地接近最优。

MDE 的一个局限性在于它假设领域之间的交互效应
可以被线性近似，即混合数据的效果
等于各领域独立效果的加权和。
但在实践中，某些领域之间存在\textbf{非线性协同}
（如代码数据和数学数据的组合效果
超过两者独立贡献之和），
或\textbf{竞争效应}
（如过多的对话数据可能降低模型的正式写作能力）。
这些交互效应在 MDE 的框架中尚未被充分建模，
是配比优化研究的一个重要开放问题。

\subsubsection{Frontier 模型的数据策略}

虽然 frontier 模型（GPT-4o、Claude 3.5、Gemini 2.0、
DeepSeek-V3、Qwen3、LLaMA 4）的详细数据配比通常不公开，
但从技术报告和社区分析中可以推断一些趋势：

\begin{itemize}[noitemsep]
  \item \textbf{代码占比持续增加}：
        从 LLaMA 1 的 $\sim$5\% 到 LLaMA 3 的 $\sim$25\%
        再到 DeepSeek-V3 的 $\sim$30\%。
        LLaMA 4 在 40T token 的训练数据中
        据推测代码占比进一步提高。
  \item \textbf{多阶段配比调整}：
        几乎所有 frontier 模型都在训练后期（annealing 阶段）
        提高高质量数据的比例。
        这已经从「实验性策略」变成了「标准做法」。
  \item \textbf{数学/推理数据成为独立类别}：
        不再被归入「学术论文」，
        而是作为独立的数据源（包含大量合成数学数据）
        在配比中占据 5\%--15\% 的份额。
\end{itemize}

\subsection{课程学习}

另一种思路是\textbf{课程学习}（Curriculum Learning）：
不是在整个训练过程中使用固定的数据配比，
而是在不同训练阶段使用不同的数据。

直觉类比：人类学习也是从简单到复杂：
先学字母，再学单词，然后学句子，最后读文章。
课程学习将类似的思想应用到 LLM 训练中。

一种常见的策略是：
\begin{enumerate}[noitemsep]
  \item \textbf{早期阶段}（前 20\% 的训练）：
        使用高质量、语言简单、结构清晰的数据
        （如 Wikipedia、教科书）。
  \item \textbf{中期阶段}（20\%--80\%）：
        逐步引入更多样化的数据
        （网页、代码、论文）。
  \item \textbf{晚期阶段}（最后 20\%，称为 annealing）：
        提高高质量数据的比例，
        同时降低学习率：
        让模型在收敛前最后一次「复习」最重要的知识。
\end{enumerate}

LLaMA 3~\cite{dubey2024llama3} 的训练
就使用了这种 annealing 策略：
在训练的最后阶段，
将高质量数据（教育内容、代码、数学）的比例提高，
并将学习率线性衰减到零。
这个简单的策略带来了显著的质量提升。

\subsubsection{Annealing 的定量细节}

Annealing（退火）已经从“实验性技巧”变成了
2025 年的\textbf{标准做法}，
但其最优配置仍然是一个活跃的研究问题。
以 LLaMA 3 为例，技术报告披露的 annealing 细节如下：
\begin{itemize}[noitemsep]
  \item \textbf{启动时机}：在 15T token 训练的最后 8\%（约 1.2T token）处开始。
  \item \textbf{学习率}：从峰值 $1.5 \times 10^{-4}$
        线性衰减到 $1.5 \times 10^{-5}$（峰值的 10\%）。
  \item \textbf{数据配比变化}：高质量教育数据从 $\sim$5\% 提升到 $\sim$20\%；
        代码数据从 $\sim$25\% 提升到 $\sim$35\%；
        数学数据从 $\sim$5\% 提升到 $\sim$15\%。
        相应地，普通网页数据的占比大幅降低。
\end{itemize}

Annealing 的效果可以从\textbf{基准测试的突变}中清楚看到。
LLaMA 3 团队报告，在 annealing 阶段开始后，
GSM8K（小学数学推理）的准确率在短短 0.2T token 内
从 65\% 跃升到 \textbf{78\%}：
这 13 个百分点的提升相当于
此前整个训练过程中最后 3T token 的累积提升量。
MATH（高中及竞赛数学）上的提升更加戏剧性：
从 28\% 跃升到 38\%。
这些数据证明了：
\textbf{在模型已经充分学习了通用语言能力之后，
在收尾阶段集中“灌注”高质量领域数据，
可以实现远超线性预期的能力提升}。

这一现象的直觉解释是：
预训练的前 92\% 让模型建立了强大的通用语言理解和推理基础：
它已经“知道”了数学的基本概念和推理模式，
但这些知识以一种\textbf{潜在的}、未被充分激活的形式存在。
Annealing 阶段的高质量数据
像一把钥匙，\textbf{激活}了这些潜在能力。
这与人类学习中的“关键期”（critical period）概念有趣地呼应：
在基础能力已经成熟之后，
针对性的强化训练可以事半功倍。

\subsubsection{课程学习的前沿方向}

课程学习的思想正在从简单的“分阶段调整配比”
演化为更精细的\textbf{逐样本难度调度}。
Albalak 等人（2024）% NEW REF: albalak2024curriculum = A Survey on Data Selection for Language Models, Albalak et al., TMLR 2024
将数据调度策略分为三个层次：
\begin{enumerate}[noitemsep]
  \item \textbf{静态配比}：在整个训练中使用固定比例（最简单，如 LLaMA 1）。
  \item \textbf{阶段性调整}：在不同训练阶段切换配比
        （如 LLaMA 3 的 annealing）。
  \item \textbf{在线调度}：基于模型\textbf{当前状态}
        动态选择每个 batch 的数据：
        类似于主动学习（active learning），
        每步都选择对当前模型“最有教育意义”的数据。
\end{enumerate}

第三层次的“在线数据选择”是 2025 年的前沿方向。
\textbf{Skill-It}（CMU + Meta，2024）% NEW REF: chen2024skillit = Skill-It: A Data-Driven Skills Framework for Understanding and Training Language Models, Chen et al., NeurIPS 2024
将训练数据组织为\textbf{技能图}（skill graph）：
每种“技能”（如语法、事实知识、逻辑推理）
对应一组训练数据，
技能之间存在依赖关系（如“逻辑推理”依赖于“语法”）。
训练调度器按照技能图的拓扑顺序排列数据：
先让模型掌握基础技能，
再逐步引入依赖于基础技能的高级技能。
在 GPT-2 规模的实验中，
Skill-It 的训练效率是随机调度的 \textbf{2 倍}：
即达到相同性能只需一半的训练步数。

\subsection{数据重复（Multi-Epoch Training）}

关于数据重复（multi-epoch training），
经验表明高质量数据重复 2--4 次是安全的，
但超过 4 次后过拟合风险急剧上升。
一个实用的启发式规则是：
如果某类数据的质量远高于平均水平
（如精选的教科书文本或高质量代码），
可以重复使用更多次；
而低质量数据（如未过滤的网页文本）
即使只用一次也要谨慎：
因为模型可能记住其中的噪声模式。

Muennighoff 等人（2023）的研究更精确地量化了这个问题：
对于 C4 数据集上的实验，
\textbf{4 epochs 的退化等价于将数据集大小减半}。
也就是说，重复使用数据 4 次
比首次使用新数据的价值低了一半。
这为「应该继续爬取新数据还是重复使用现有数据」
提供了定量的决策依据。
% NEW REF: muennighoff2023scaling = Muennighoff, Niklas and others. Scaling Data-Constrained Language Models. In NeurIPS, 2023.

\section{合成数据：模型教模型}

2024--2026 年的一个重大趋势是\textbf{合成数据}：
用已有的强大模型生成训练数据来训练新模型。

\subsection{教科书级合成数据：Phi 的启示}

微软的 Phi 系列模型是合成数据成功应用的标杆。
Phi-1（2023 年）使用了「教科书级」
（textbook quality）的合成数据
来训练一个仅 1.3B 参数的代码生成模型。
关键洞察是：与其在大量的低质量 GitHub 代码上训练，
不如让 GPT-3.5 生成
\textbf{像教科书一样清晰、系统、循序渐进}的代码教程。
% NEW REF: gunasekar2023phi = Gunasekar, Suriya and others. Textbooks Are All You Need. arXiv preprint arXiv:2306.11644, 2023.

Phi-1 仅用约 7B token 的训练数据
（其中大部分是合成的），
在 HumanEval 代码基准测试上
达到了与 StarCoder（15B 参数、1T token 训练）
相当的水平：
参数量少 10 倍，数据量少 100 倍。
这个结果震撼了整个社区。

微软的 Phi-4 将这一策略推向了极致：
Phi-4 在\textbf{预训练阶段}就大量使用了合成数据：
其预训练语料中包含大量由 GPT-4 生成的高质量合成内容，
涵盖推理链、教科书式讲解和结构化问答等多种形式。
在后续的 SFT 阶段，
它使用了约 140 万个精心策划的样本进行监督微调。
结果是一个仅 14B 参数的模型，
在多项任务上超越了比它大数倍的模型：
证明了\textbf{数据质量远比数据数量重要}。

Phi 系列的成功揭示了“教科书级”合成数据的三个设计原则：

\begin{enumerate}[noitemsep]
  \item \textbf{循序渐进}：
        每段合成文本都从简单概念出发，
        逐步引入更复杂的概念：
        像教科书而非百科全书。
        例如，一段关于递归的合成代码教程
        会先展示阶乘函数，
        然后是斐波那契数列，
        最后是树遍历：
        而不是随机罗列这三个例子。
  \item \textbf{自包含}：
        每段文本应该是自包含的：
        读者不需要阅读其他材料就能理解。
        这与 Wikipedia 的风格形成对比：
        Wikipedia 依赖大量的交叉引用，
        而教科书级合成数据在单段内解释所有必要的背景。
  \item \textbf{练习与解答}：
        每个概念后面跟着练习题和详细解答。
        这种“教学—练习—反馈”的循环
        提供了比纯文本更丰富的学习信号：
        模型不仅要理解概念，
        还要学会\textbf{应用}概念。
\end{enumerate}

从 Phi-1 到 Phi-4 的演进轨迹清晰地展示了
合成数据策略的\textbf{规模化路径}：

\begin{center}\small
\renewcommand{\arraystretch}{1.3}
\begin{tabular}{lrrll}
\toprule
\rowcolor{figLightGray}
\textbf{模型} & \textbf{参数} & \textbf{训练数据} & \textbf{合成数据策略} & \textbf{亮点} \\
\midrule
Phi-1 (2023) & 1.3B & $\sim$7B tok & GPT-3.5 生成代码教程 & 代码任务媲美 15B \\
Phi-1.5 (2023) & 1.3B & $\sim$30B tok & 教科书 + 推理链 & 通用推理 \\
Phi-2 (2023) & 2.7B & $\sim$250B tok & 大规模多领域合成 & 部分超越 25B \\
Phi-3 (2024) & 3.8B & 4.9T tok & 合成 + 过滤网页 & 接近 GPT-3.5 \\
Phi-4 (2024) & 14B & 9.8T tok & 预训练大量合成 & 超越 GPT-4-turbo(部分) \\
\bottomrule
\end{tabular}
\end{center}

注意一个关键的转变：
Phi-1/1.5 的核心论点是“少量精选数据胜过大量数据”，
但到 Phi-3/4 时，训练数据量已经增长到 T 级别：
只是在这个规模下，合成数据占了很大比例。
这说明“教科书级数据”的优势
不是因为数据量\textbf{少}，
而是因为质量\textbf{高}：
当你能大规模生产高质量合成数据时，
规模仍然是有益的。

\subsection{自指令与进化指令}

\textbf{Self-Instruct}（Wang 等人，2023）
是合成数据生成的另一种重要范式。
它的思路是让 LLM 自己生成指令数据：
给模型少量种子指令（如 175 条），
让模型根据这些种子生成新的指令和对应的回答，
然后用这些数据来微调模型本身。
% NEW REF: wang2023selfinstruct = Wang, Yizhong and others. Self-Instruct: Aligning Language Models with Self-Generated Instructions. In ACL, 2023.

Self-Instruct 的流水线包含四个步骤：
\begin{enumerate}[noitemsep]
  \item \textbf{指令生成}：
        从种子指令池中随机抽取 8 条作为 in-context 示例，
        让 GPT-3 生成新的指令。
  \item \textbf{指令分类}：
        判断新指令是“分类任务”还是“生成任务”：
        这决定了后续生成回答的方式。
  \item \textbf{实例生成}：
        对于分类任务，先生成输出（标签），
        再基于标签反向生成输入（避免标签偏差）；
        对于生成任务，先生成输入再生成输出。
  \item \textbf{质量过滤}：
        移除与种子指令过于相似的（ROUGE-L $>$ 0.7）、
        过长或过短的、以及包含“As an AI model”等
        元认知表述的低质量样本。
\end{enumerate}

这个看似简单的流水线
从 175 条种子指令出发，
产生了 \textbf{52,000 条}高质量指令：
300 倍的放大率。
更令人惊讶的是，
使用这些自动生成的数据微调后的 GPT-3
在 SuperNI（一个包含 756 种 NLP 任务的基准测试）上
与使用人工标注数据微调的 InstructGPT
\textbf{仅相差 5\%}：
证明了在一定质量范围内，
LLM 自我生成的训练数据
可以替代昂贵的人工标注。

\textbf{Evol-Instruct}（WizardLM 团队）
在 Self-Instruct 的基础上更进一步，
通过\textbf{进化}策略来增加指令的复杂度。
具体做法是：
从简单的指令出发，
让 LLM 对指令进行“深化”（添加约束）、
“拓宽”（增加子任务）、
“具体化”（添加背景）等进化操作，
生成更复杂、更具挑战性的指令。
多轮进化后，生成的指令
从简单的“翻译这段话”
变成了“将这段 19 世纪英文法律文件翻译成现代中文，
保留原始的法律术语，并在注释中解释关键概念”。
% NEW REF: xu2023wizardlm = Xu, Can and others. WizardLM: Empowering Large Language Models to Follow Complex Instructions. arXiv preprint arXiv:2304.12244, 2023.

Evol-Instruct 的每轮进化有五种“变异算子”：
(1) \textbf{添加约束}——“写一个排序算法” $\to$ “写一个时间复杂度 $O(n\log n)$ 的排序算法，不能使用额外空间”；
(2) \textbf{增加推理深度}——“计算 $3 + 5$” $\to$ “证明 $\sum_{k=1}^{n} k = n(n+1)/2$”；
(3) \textbf{具体化}——“写一段代码” $\to$ “用 Python 实现一个 REST API，使用 FastAPI 框架，支持 CRUD 操作”；
(4) \textbf{增加输入复杂度}——扩展输入数据的规模或复杂度；
(5) \textbf{替换为更罕见的要求}，用不常见的概念替换常见概念。
经过 4 轮进化，指令的平均“难度分数”（由 GPT-4 评估）
从 2.3/10 提升到 \textbf{7.1/10}：
覆盖了从简单到高难度的完整难度谱。

Self-Instruct 和 Evol-Instruct 的\textbf{质量--多样性权衡}
是合成数据领域的核心挑战之一。
使用更强的模型（如 GPT-4 而非 GPT-3.5）
生成的指令质量更高，但\textbf{风格更均匀}：
GPT-4 倾向于生成结构化、条理清晰的回答，
这使得模型在微调后也倾向于这种风格，
失去了回答风格的多样性。
一种缓解策略是\textbf{模型多样性}：
使用多个不同的生成模型
（GPT-4、Claude、Mixtral、LLaMA）
各自生成一部分指令数据，
然后混合使用。
不同模型的“风格签名”不同：
GPT-4 倾向于结构化的要点列表，
Claude 偏好流畅的长段落叙述，
LLaMA 的风格更多变：
混合后的数据集在风格多样性上
远优于单一模型生成。

\subsection{大规模合成预训练数据}

合成数据的应用已经从 SFT（微调）阶段
扩展到了\textbf{预训练}阶段——这是一个质的变化。

\textbf{Cosmopedia}（HuggingFace，2024 年）
是目前最大的开源合成预训练数据集。
它使用 Mixtral-8x7B-Instruct 生成了
\textbf{超过 3000 万个文件、250 亿 token} 的合成内容，
包括教科书、博客文章、故事和 WikiHow 式文章。
Cosmopedia 的关键设计是\textbf{提示多样性}：
使用了数百个不同的主题和写作风格模板，
确保合成内容的多样性，避免模型坍缩的风险。
Cosmopedia 被用于训练 SmolLM 系列小模型，
在多项基准上取得了与使用真实数据训练的模型
相当甚至更好的结果。
% NEW REF: benallal2024cosmopedia = Ben Allal, Loubna and Lozhkov, Anton and others. Cosmopedia: how to create large-scale synthetic data for pre-training. HuggingFace Blog, 2024.

Cosmopedia 的\textbf{提示工程}
是其成功的核心：
而不仅仅是调用 Mixtral 那么简单。
HuggingFace 团队设计了 \textbf{33 种主题类别}
和 \textbf{8 种写作风格}，
交叉组合产生了超过 264 种不同的提示模板。
每种模板规定了：目标读者的年龄/知识水平、
文章的类型（教程/叙事/对话/百科）、
期望的长度和复杂度。
例如，一个模板可能要求：
“为一个 10 岁的孩子写一篇关于光合作用的故事，
使用简单的类比和生动的描述，长度约 500 词”。
另一个模板则要求：
“为研究生撰写一篇关于光合作用的技术综述，
引用关键实验结果和化学方程式，长度约 2000 词”。
同一个主题在不同模板下产出的文本
在词汇、句法和概念深度上完全不同：
这正是多样性的来源。

为了防止生成数据的\textbf{模式坍缩}，
Cosmopedia 还引入了“种子文本”机制：
从 Wikipedia、Wikidata 和 Reddit 帖子中
随机抽取一段真实文本作为“种子”，
要求模型基于种子的\textbf{主题}（但非内容）
来生成新文本。
种子的多样性确保了生成主题的广泛覆盖，
而“基于主题而非复述内容”的要求
避免了简单的改写式生成。
这种设计使得 Cosmopedia 的\textbf{主题分布}
与 Wikipedia 的主题分布高度对齐：
这在一定程度上缓解了合成数据
通常过度集中于热门话题的倾向。

\textbf{Cosmopedia v2}（2024 年底）进一步改进了生成策略：
引入了\textbf{多轮改写}（multi-turn rewriting）：
先让模型生成初稿，
然后让另一个模型（或同一模型的不同提示）
对初稿进行质量评分和改写，
最终只保留得分最高的版本。
这种“生成-评估-改写”的循环
使 v2 的每 token 平均教育价值（以下游基准衡量）
比 v1 提高了约 15\%。

\textbf{OpenHermes 2.5}（Teknium，2024 年）
是开源社区最成功的合成 SFT 数据集之一。
它从 15 个不同的数据源中策划了
\textbf{超过 100 万条}高质量合成指令和对话样本：
包括 GPT-4 生成的数据、数学推理数据（MetaMath）、
代码数据（Glaive Code Assistant）和推理链数据（CoT Alpaca）等。
OpenHermes 2.5 被用于训练了多个在社区排行榜上名列前茅的模型，
证明了\textbf{精心策划的多源合成数据}
可以在微调阶段带来显著的质量提升。
% NEW REF: teknium2024openhermes = Teknium. OpenHermes 2.5 Dataset. HuggingFace, 2024.

\textbf{NVIDIA Nemotron-4 340B}（2024 年）
代表了合成数据策略的另一个极端：
它用 \textbf{98\% 的合成数据}训练了一个 340B 参数的模型，
在多项基准上与 GPT-4 竞争。
NVIDIA 开放了完整的合成数据生成管线，
允许开发者使用 Nemotron-4 家族模型
（包含生成器、评分器和奖励模型）
为自己的应用领域生成定制化的合成数据。
这标志着合成数据从「研究技巧」走向了「工业基础设施」。
% NEW REF: nvidia2024nemotron4 = Adler, Bo and others. Nemotron-4 340B Technical Report. arXiv preprint arXiv:2406.11704, 2024.

\subsection{数学与代码的合成数据}

合成数据在数学和代码领域特别有效，
因为这些领域有\textbf{可验证的正确性标准}。

\textbf{数学合成数据}：
可以自动生成数学题目和解答过程，
然后通过符号计算（如 SymPy）验证答案的正确性。
DeepSeekMath~\cite{shao2024grpo} 就利用了
大量自动生成和验证的数学数据
来提升模型的数学推理能力。

\textbf{MetaMathQA}（Yu 等人，2024）% NEW REF: yu2024metamathqa = MetaMath: Bootstrap Your Own Mathematical Questions for Large Language Models, Yu et al., ICLR 2024
是数学合成数据的标杆工作。
其核心策略是\textbf{数学题目的自举}（bootstrapping）：
从 GSM8K 和 MATH 的原始题目出发，
使用 GPT-3.5-Turbo 对题目进行多种变换：
\textbf{答案增强}（让模型生成更详细的解题步骤）、
\textbf{题目重述}（用不同的语言重新描述同一道题）、
\textbf{FOBAR}（从答案反推出题目中的条件）、
以及\textbf{Self-Verification}
（让模型检查自己的解答是否正确）。
这些变换将 GSM8K 的 7,473 道原始题目
扩展为 \textbf{395K} 道变体题目：
每道原始题目平均产出 53 道不同表述和难度的变体。

MetaMathQA 的一个关键发现是
\textbf{答案增强}对性能提升贡献最大：
将简短的答案（“答案是 42”）
扩展为详细的分步推理链
（“首先，根据条件 A 可得 X...
其次，将 X 代入等式 B...
因此答案是 42”）
使模型从“记忆答案”转变为“学习解题过程”。
使用 MetaMathQA 微调的 7B 模型
在 GSM8K 上达到了 \textbf{66.5\%} 的准确率：
比使用原始 GSM8K 数据微调的同规模模型高出 20 个百分点。

\textbf{NuminaMath}（Numina 团队，2024）
是另一个重要的数学合成数据集。
它收集了来自数学竞赛（AMC、AIME、IMO）、
在线判题系统（AoPS、Project Euler）
和数学教材的 \textbf{超过 86 万道题目和解答}。
对于没有现成解答的竞赛题，
NuminaMath 使用 GPT-4 生成候选解答，
然后通过\textbf{符号验证}（SymPy 数值检查 + 形式化验证器）
来过滤错误答案。
这种“生成-验证”的循环使得合成解答的正确率
从初始的约 60\% 提升到验证后的 95\%+。
% NEW REF: numina2024 = NuminaMath: A Dataset for Mathematical Reasoning, Numina Team, 2024.

\textbf{代码合成数据}：
类似地，合成的代码可以通过\textbf{执行测试}来验证：
生成代码和对应的单元测试，
只保留能通过测试的代码作为训练数据。
这种“执行验证”（execution-based verification）
提供了比人工标注更精确的质量信号。

代码合成数据的\textbf{质量--多样性权衡}
比文本数据更加微妙。
高质量的代码要求语法正确、逻辑正确、
遵循编码规范、处理边界情况：
而这些维度之间可能存在冲突。
例如，一个“正确但晦涩”的解法
和一个“清晰但效率不佳”的解法，
哪个对训练更有价值？
实践中的共识是：
在\textbf{预训练}阶段，多样性优先：
各种风格和质量等级的代码都有价值；
在\textbf{SFT}阶段，质量优先：
只保留“教科书级”的清晰、正确的代码。
这种分阶段策略在 DeepSeek-Coder 和 StarCoder 2
的训练中都被采用。

\subsection{合成数据在 frontier 模型中的占比}

截至 2025--2026 年，合成数据在 frontier 模型训练中的占比
已经从早期实验的「少量补充」增长到了\textbf{实质性组成部分}：

\begin{itemize}[noitemsep]
  \item \textbf{预训练阶段}：Phi-4 和 Nemotron-4 表明
        合成数据可以占预训练语料的 30\%--98\%：
        具体比例取决于合成数据的质量和多样性。
        但对于 frontier 级模型（如 GPT-4o、Claude 3.5），
        普遍估计合成数据占比在 20\%--40\% 之间。
  \item \textbf{SFT 阶段}：合成数据已经成为主流。
        几乎所有 frontier 模型的 SFT 数据中，
        合成指令和合成对话占比超过 50\%。
  \item \textbf{RLHF/DPO 阶段}：
        合成偏好数据（如 Constitutional AI 生成的偏好对）
        正在逐步取代成本高昂的人类标注。
\end{itemize}

这个趋势的驱动力是清晰的：
高质量的真实数据正在\textbf{枯竭}。
互联网上可用的高质量文本是有限的，
而模型对数据的需求仍在增长。
合成数据是解决这一矛盾的关键路径：
但它引入的模型坍缩风险也必须被认真对待。

\subsection{合成数据的风险：模型坍缩}

合成数据并非万能药。
Meta FAIR 训练了超过 1000 个 LLM
来研究合成数据的最优使用方式%
\footnote{该结论来自 Meta FAIR 的系统性实验，具体论文待补充。}。
核心结论是：预训练的最优比例约为 30\% 合成 + 70\% 真实数据，
可以加速训练 5--10 倍。

但需要警惕\textbf{模型坍缩}（model collapse）。
Shumailov 等人（2023）的研究发现了一个令人不安的现象：
如果模型 A 生成的数据被用来训练模型 B，
B 生成的数据再用来训练模型 C，如此递归数代，
模型的输出会逐渐\textbf{退化}：
「遗忘」数据分布的尾部（低频但重要的模式），
输出越来越同质化、越来越「平庸」。
% NEW REF: shumailov2023collapse = Shumailov, Ilia and others. The Curse of Recursion: Training on Generated Data Makes Models Forget. arXiv preprint arXiv:2305.17493, 2023.

这个现象的直觉解释是：
模型在生成数据时会\textbf{放大高频模式、忽略低频模式}：
类似于反复复印一张照片，每次都会损失一些细节。
经过多代递归后，只有最高频的模式被保留，
丰富多样的尾部分布完全消失。

\subsubsection{模型坍缩的两种形式}

Shumailov 等人将模型坍缩区分为两个阶段：

\textbf{早期坍缩}（early collapse）：
模型首先丢失分布尾部的信息：
低频但有意义的模式（如罕见的写作风格、
小众领域的知识、少数群体的表达方式）逐渐消失。
这个阶段的坍缩往往不易察觉，
因为主流基准测试通常不测量尾部分布的保留程度。

\textbf{晚期坍缩}（late collapse）：
经过多代递归后，模型输出的分布
与原始数据分布完全不同：
方差急剧减小，输出变得高度同质化。
在极端情况下，模型可能退化为
反复生成少数几种固定模式。

\subsubsection{2025 年的缓解策略}

2025 年的研究提出了多种缓解模型坍缩的新方法：

\begin{itemize}[noitemsep]
  \item \textbf{合成数据检测与过滤}：
        在训练数据中检测并移除由 LLM 生成的文本。
        随着互联网上 AI 生成内容的比例急剧上升
        （估计已超过某些领域网页内容的 10\%--20\%），
        即使不主动使用合成数据，
        从 Common Crawl 中爬取的网页也可能包含大量 AI 生成文本。
        这使得合成数据检测成为数据清洗流水线的新必备步骤。
  \item \textbf{多样性正则化}：
        在合成数据生成时显式优化输出的多样性：
        例如使用高温度采样、多样化的提示模板、
        以及刻意覆盖长尾主题和风格。
  \item \textbf{真实数据锚定}：
        始终保留一定比例的真实数据作为「锚」，
        防止分布漂移。Nemotron-CC 的「改写」策略
        是一个优雅的折中：
        使用 LLM 改善真实数据的表达质量，
        同时保留真实数据的内容多样性。
  \item \textbf{分布匹配验证}：
        定期检查合成数据的分布
        是否偏离真实数据的分布，
        使用统计检验（如 KL 散度、MMD）
        来量化偏离程度并及时干预。
\end{itemize}

\begin{warnbox}[避免模型坍缩的实践建议]
(1) \textbf{始终混入真实数据}：
合成数据应该是真实数据的\textbf{补充}而非\textbf{替代}。
(2) \textbf{限制递归深度}：
不要用模型 A 生成的数据训练模型 B，
再用 B 生成的数据训练 C。
尽量始终用最强的模型（如 GPT-4）作为生成源。
(3) \textbf{多样性保障}：
在合成数据中刻意注入多样性
（不同的提示模板、不同的主题、不同的风格）。
(4) \textbf{质量过滤}：
不是所有合成数据都是好的。
使用验证机制（执行测试、一致性检查）
过滤掉低质量的合成样本。
(5) \textbf{检测互联网上的 AI 生成内容}：
2025 年后从 Common Crawl 爬取的数据
可能已经包含大量 AI 生成文本。
不加检测地使用这些数据会导致隐性的模型坍缩。
\end{warnbox}

\subsection{合成数据的不可或缺时代（2026）}

2026 年正在成为合成数据从「可选的增强手段」
转变为「不可或缺的基础设施」的分水岭。
驱动这一转变的核心原因是：
\textbf{互联网上的高质量自然文本正在逼近枯竭}。
当 Qwen3-Max 的训练数据达到 \textbf{36T token}、
GLM-5 使用 \textbf{28.5T token} 时：
相比 GPT-3 时代的 300B token，
数据规模已经膨胀了两个数量级：
依靠纯自然数据已经难以为继。

\subsubsection{模型坍缩的实证与应对}

模型坍缩不再是理论警告，而是\textbf{经过实证验证的威胁}。
2025--2026 年的多项研究系统性地证实了
在 AI 生成数据上递归训练会导致分布退化。
关键在于建立可预测的合成数据使用框架，
而非简单地避免使用合成数据。

Dohmatob 等人（2024，arXiv:2402.01942）从信息论的角度
给出了模型坍缩的理论分析框架。
他们证明了在高斯混合模型的简化设定下，
递归训练 $n$ 代后，估计误差以
$O(\sqrt{n / N})$ 的速率增长
（$N$ 为每一代的训练样本数）。
当合成数据占比超过某个临界值时，
模型的测试误差开始\textbf{单调递增}：
即使增加总训练数据量也无法逆转。
这个临界值取决于合成数据的质量和原始数据的比例：
对于高质量合成数据（如 GPT-4 生成），
安全阈值约为 40--60\%；
对于低质量合成数据（如小模型生成），
阈值可能低至 10--20\%。
% NEW REF: dohmatob2024collapse = Dohmatob, Elvis and others. A Tale of Tails: Model Collapse as a Change of Scaling Laws. arXiv preprint arXiv:2402.01942, 2024.

\textbf{坍缩的定量检测方法}也在 2025 年逐渐成熟。
Guo 等人（2024，arXiv:2404.10040）提出了基于
\textbf{token 频率分布}的坍缩检测指标：
在正常训练中，模型生成文本的 token 频率
应该近似服从 Zipf 定律；
当坍缩发生时，频率分布的尾部会急剧收缩，
表现为 Zipf 系数从 $\sim$1.0 偏移到 $>$1.5。
这种简单的统计检验可以作为训练过程中的
实时“哨兵”，在坍缩的早期阶段就发出警报。
另一种方法是监控模型输出的
\textbf{Self-BLEU 分数}（生成多个回答之间的 BLEU 相似度）：
Self-BLEU 持续上升表明模型的输出多样性正在降低。
% NEW REF: guo2024collapse = Guo, Rui and others. Curious Decline of Linguistic Diversity: A Large-Scale Analysis of LLM-Generated Text. arXiv preprint arXiv:2404.10040, 2024.

在\textbf{缓解策略}方面，
除了前文提到的“始终混入真实数据”的经验法则，
更精细的方法正在出现。
\textbf{数据来源加权}策略根据数据的“代际”
（真实数据 = 第 0 代，由第 0 代模型生成 = 第 1 代，以此类推）
分配递减的权重——第 $k$ 代数据的权重为 $w_k = \alpha^k$
（$\alpha < 1$，通常取 0.5--0.8），
这样递归生成的数据在训练目标中的影响被指数级压低。
\textbf{对比训练}（contrastive training）
则显式地训练模型区分真实数据和合成数据的分布特征，
在损失函数中加入一个惩罚项，
鼓励模型维持对真实数据分布的保真度。
这些方法的核心思想是一致的：
\textbf{合成数据的价值在于补充真实数据的不足，
而非替代真实数据的统计特性}。

\subsubsection{合成数据 Scaling Laws}

COLM 2025 发表、2026 年 3 月更新的研究
（arXiv:2503.19551v2）首次建立了
\textbf{合成数据的 Scaling Laws}：
证明合成数据在预训练中的使用具有
\textbf{可预测的规模化规律}。
这意味着团队可以在小规模实验中
准确预估大规模使用合成数据时的效果，
而无需盲目试错。
这一发现从根本上改变了合成数据的工程实践：
从「经验性的混合比例调整」
走向「基于 Scaling Laws 的系统性规划」。

\subsubsection{工业级合成数据管线}

微软的 \textbf{SynthLLM} 代表了
合成数据从实验室走向工业级基础设施的标志。
SynthLLM 提供了一套可扩展的合成数据生成管线，
目标是突破所谓的「AI 数据墙」：
即高质量自然数据的供给上限。
其核心设计包括多阶段生成、
自动质量验证和多样性保障机制，
使得大规模合成数据的生产变得可靠和可重复。

SynthLLM 的管线架构包含三个关键阶段。
\textbf{生成阶段}使用多个不同的生成模型
（而非单一模型）并行生成合成数据，
利用模型多样性来减少单一模型偏差导致的分布坍缩。
\textbf{验证阶段}对生成的数据进行多维度质量检查：
包括事实一致性验证（通过外部知识库交叉检查）、
逻辑连贯性检测（通过 NLI 模型评估段落间的蕴含关系）、
以及重复性检测（确保合成数据不会过度集中在少数模式上）。
\textbf{多样性保障阶段}通过主动采样策略
确保合成数据覆盖长尾主题和罕见知识领域：
这一阶段使用嵌入空间的覆盖度指标
来量化当前合成数据相对于目标分布的覆盖率，
并优先生成低覆盖区域的数据。
这种“精心练习”（deliberate practice）的理念：
聚焦模型的薄弱环节而非均匀生成：
是 SynthLLM 区别于简单的“大量生成 + 过滤”策略的核心。

\subsubsection{数据增强的新路径}

除了直接合成，
前沿团队也在探索\textbf{数据增强}的新路径。
据报道，DeepSeek 与百度在 AI 搜索数据方面展开合作，
通过搜索引擎获取的结构化查询--回答对
作为训练数据的补充来源。
这种「搜索数据增强」策略
利用了搜索引擎对用户意图的理解，
能产生更贴近真实用户需求的训练数据。

\textbf{回译增强}（backtranslation augmentation）
是另一种经过验证的数据增强路径。
将英文文档翻译成中文（或其他目标语言），
再翻译回英文，可以生成语义等价但表述不同的文本：
这种“同义改写”增加了数据的表面多样性，
而不引入语义噪声。
Google 在多语言模型训练中广泛使用了这一策略，
通过 T5 翻译模型对低资源语言的数据进行增强。

\textbf{知识图谱驱动的合成}是更结构化的方向。
从 Wikidata 等知识图谱中提取实体关系三元组，
然后让 LLM 将这些结构化知识“展开”为自然语言文本。
这种方法的优势是生成的内容具有可验证的事实准确性
（因为原始知识来自结构化数据库），
但生成的文本风格往往过于百科全书化。
结合搜索数据增强和知识图谱驱动，
可以在保证事实准确性的同时
兼顾文本风格的多样性：
这是 2026 年数据增强研究的一个活跃方向。

\begin{warnbox}[2026 年的合成数据现实]
合成数据已经不是一个选择题。
当训练数据需求达到数十万亿 token 级别时，
任何 frontier 模型都无法仅依赖自然数据。
关键的工程问题已从「是否使用合成数据」
转变为「如何安全、可预测地使用合成数据」：
Scaling Laws 的建立、
模型坍缩检测机制的成熟、
以及工业级生成管线的出现，
正在让这一转变变得可行。
\end{warnbox}

\subsection{Constitutional AI 与 AI 反馈}

Anthropic 的 Constitutional AI~\cite{bai2022constitutional}
开创了另一种合成数据的使用方式：
用 AI 生成\textbf{反馈}而非\textbf{内容}。

传统的 RLHF 需要人类标注者
对模型的输出进行偏好排序：
这个过程既昂贵又难以规模化。
Constitutional AI 用一组明确的「原则」
（宪法，constitution）来指导 AI 自我评估：
模型先生成回答，
然后根据宪法中的原则对自己的回答进行批评和修订，
最后用修订后的版本作为偏好学习的训练数据。

这种方法极大地降低了获取偏好数据的成本，
使得快速迭代 RLHF 训练成为可能。

\section{数据处理工具生态}

2025--2026 年，LLM 的数据处理已经形成了成熟的工具生态。

\subsection{Datatrove}

\textbf{Datatrove}（HuggingFace 开源）是 FineWeb 背后的数据处理框架。
它将整个流水线抽象为「读取 $\to$ 过滤 $\to$ 去重 $\to$ 写入」的模块化流程，
支持在单机和集群上运行。
核心设计理念是\textbf{可复现性}：
每一步的参数和随机种子都被记录，
确保相同的输入总是产生相同的输出。

Datatrove 的架构设计值得学习：
每个处理步骤（称为 “block”）
都实现统一的接口（读取文档流、输出文档流），
可以自由组合。
这种设计使得实验不同的过滤策略变得非常容易：
只需替换或添加 block 即可。

\subsection{CCNet}

\textbf{CCNet}（Meta 开源）
是专门为 Common Crawl 处理设计的工具。
它实现了完整的清洗流水线：
语言识别（fastText）$\to$
段落去重（SHA-256 哈希）$\to$
困惑度过滤（KenLM）。
LLaMA~\cite{touvron2023llama} 的训练数据
就是通过 CCNet 处理的。

CCNet 的设计哲学是\textbf{简洁和可扩展}：
整个流水线由三个解耦的阶段组成，
每个阶段可以独立运行和并行化。
困惑度过滤使用 5-gram KenLM 模型，
在 Wikipedia 上训练，
将文档分为“头部”（困惑度 $< 230$，约占 33\%）、
“中部”和“尾部”三个质量等级。
LLaMA 系列选择了“头部”文档作为预训练数据：
这一简单的策略在 2023 年被证明足以
构建有竞争力的训练数据。
CCNet 的另一个重要贡献是\textbf{分片处理}机制：
它将 Common Crawl 按 URL 的哈希值分片，
使得去重操作可以在分片内独立完成，
大幅降低了内存需求。

\subsection{Data-Juicer}

\textbf{Data-Juicer}（阿里巴巴 DAIL 开源）
提供了 50+ 种数据处理算子（operator），
涵盖文本质量评估、语言识别、数据多样性分析、
安全内容过滤等。它的特色是
支持「数据菜谱」（data recipe）的概念：
将一组处理步骤和参数打包为可复用的配置文件，
便于实验和分享。

Data-Juicer 的架构遵循“算子即组件”的设计理念，
每种数据处理操作被封装为一个标准化的算子，
分为三类：\textbf{Mapper}（一对一变换，如文本清洗）、
\textbf{Filter}（一对零或一的过滤，如质量筛选）和
\textbf{Deduplicator}（多对一的去重）。
这种分类使得用户可以像搭积木一样
组合不同的处理步骤来构建定制化的数据处理流水线。
Data-Juicer 还内置了\textbf{数据分析}功能：
在处理流水线执行前后，
自动生成数据的统计分析报告
（token 长度分布、语言分布、质量分数分布等），
帮助工程师理解每一步处理的效果。
对于中文数据处理，Data-Juicer 具有独特优势：
它内置了针对中文分词、繁简转换、
中英混合文本处理等场景的专用算子：
这是 Datatrove 等以英文为中心的工具所缺乏的。
% NEW REF: chen2024datajuicer = Chen, Daoyuan and others. Data-Juicer: A One-Stop Data Processing System for Large Language Models. In SIGMOD, 2024.

\subsection{其他重要工具}

\begin{itemize}[noitemsep]
  \item \textbf{Dolma toolkit}（AI2）：
        Dolma 数据集的完整构建工具链，
        从原始数据到最终数据集的每一步都可复现。
  \item \textbf{text-dedup}：
        一个专注于文本去重的 Python 库，
        实现了 MinHash、SimHash、后缀数组等多种去重算法。
  \item \textbf{trafilatura}：
        Python 实现的网页文本提取工具，
        支持从 HTML 中精准提取正文内容。
  \item \textbf{fastText}：
        Meta 开源的高效文本分类和语言识别工具，
        在语言分类任务上是事实标准。
  \item \textbf{KenLM}：
        高效的 $n$-gram 语言模型实现，
        常用于困惑度过滤。
\end{itemize}

\subsection{数据处理流水线示例}

一个典型的数据处理流水线的伪代码如下：

\begin{lstlisting}[caption={Data processing pipeline (pseudocode)}]
from datatrove import Pipeline

pipeline = Pipeline([
    # Step 1: Read raw HTML from Common Crawl
    WARCReader(data_path="s3://commoncrawl/..."),

    # Step 2: Extract text from HTML
    TrafilaturaExtractor(),

    # Step 3: Language identification
    FastTextLanguageFilter(languages=["en", "zh"]),

    # Step 4: Rule-based quality filters
    GopherQualityFilter(
        min_doc_words=50,
        max_doc_words=100000,
        min_avg_word_length=3,
        max_symbol_to_word_ratio=0.1,
    ),

    # Step 5: Perplexity-based quality scoring
    C4QualityFilter(exclusion_writer=...),

    # Step 6: PII removal (emails, phone numbers)
    PIIRemovalFilter(),

    # Step 7: Exact deduplication
    URLDedup(),

    # Step 8: Fuzzy deduplication
    MinhashDedup(
        n_grams=5,
        num_hashes=128,
        similarity_threshold=0.8,
    ),

    # Step 9: Write processed data
    JsonlWriter(output_path="processed/"),
])

pipeline.run(num_workers=64)
\end{lstlisting}

这段伪代码展示了从原始网页到训练数据的完整流程。
每一步都是一个可配置的算子，
可以独立运行和调试。
\texttt{num\_workers=64} 表示使用 64 个并行进程处理：
处理一个 Common Crawl 快照（通常包含数十亿网页）
需要数百 CPU 核心运行数天。

\subsection{规模化的工程挑战}

在 TB 到 PB 级别处理数据，
工程挑战远超算法本身。
以 FineWeb 处理 96 个 Common Crawl 快照为例：

\begin{itemize}[noitemsep]
  \item \textbf{存储}：原始数据约 300TB，
        中间结果和最终输出需要同等甚至更多的存储空间。
  \item \textbf{计算}：全量 MinHash 去重
        需要在数百台机器上运行，
        单次去重耗时数天。
  \item \textbf{调度}：数据处理的各步骤之间存在依赖关系
        （必须先提取文本才能去重），
        需要精心设计任务调度。
  \item \textbf{容错}：在数天的运行中，
        机器故障是常态而非例外。
        流水线必须支持断点续传和部分重跑。
\end{itemize}

HuggingFace 处理 FineWeb 时使用了约 \textbf{300 万 CPU 小时}：
这是一个惊人的数字，
说明数据处理的计算成本
虽然远低于模型训练，但绝非微不足道。

\section{数据污染检测}

随着 LLM 在各种基准测试上取得越来越高的分数，
一个关键问题浮出水面：
\textbf{这些成绩到底反映了真正的泛化能力，
还是模型只是「记住」了测试数据？}

\subsection{数据污染的定义与危害}

\textbf{数据污染}（data contamination）指的是
模型的训练数据中包含了
评估基准测试的题目或答案。
由于 Common Crawl 收录了互联网上的大量内容，
而许多基准测试的数据也公开发布在网上，
训练数据与测试数据之间的重叠几乎是不可避免的。

数据污染的危害在于：
它使得基准测试的结果\textbf{不再可靠}：
我们无法区分模型是「理解了问题的模式」
还是「记住了具体的答案」。
对于推动 LLM 进步，这是一个根本性的问题：
如果基准测试失去了可信度，
我们就失去了衡量进步的尺度。

\subsection{污染的规模有多大？}

2025 年的多项研究系统性地量化了数据污染的普遍程度，
结果令人警醒。

Sainz 等人（2024）% NEW REF: sainz2024contamination = NLP Evaluation in trouble: On the Need to Measure LLM Data Contamination for each Benchmark, Sainz et al., EMNLP 2024
对 12 个主流 LLM 在 40 个基准测试上进行了污染检测。
他们发现：
在 480 个（模型 $\times$ 基准）组合中，
有 \textbf{43\%} 的组合存在可检测的数据污染。
更令人不安的是，污染程度与模型在基准上的“超常表现”
之间存在\textbf{显著的正相关}（Pearson $r = 0.67$）：
即一个模型在某个基准上的得分
越是“超出其在相似基准上的预期表现”，
越可能是因为它见过这个基准的数据。

Yang 等人（2023）% NEW REF: yang2023rethinking = Rethinking Benchmark and Contamination for Language Models with Rephrased Samples, Yang et al., arXiv 2023
发现了更隐蔽的污染形式：
即使在原始基准数据不在训练集中的情况下，
如果训练数据中包含了同一问题的\textbf{同义改写版本}
（例如来自教育网站对考试题的讨论），
模型仍然可以获得不公平的优势。
他们提出了\textbf{改写后测试}（rephrased evaluation）：
将基准测试的题目用 GPT-4 改写后重新评估。
结果显示，多个模型在改写后的版本上
得分下降了 \textbf{5--15 个百分点}：
暗示它们的部分“能力”来自记忆而非真正的理解。

\subsection{检测方法}

2024--2025 年，研究社区发展了多种污染检测方法，
从简单的字符串匹配到精妙的统计推断：

\begin{itemize}[noitemsep]
  \item \textbf{$n$-gram 重叠检测}：
        最直接的方法——检查训练数据中是否包含
        与基准测试题目或答案相同的长 $n$-gram（通常 $n \geq 13$）。
        LLaMA~\cite{touvron2023llama} 和
        GPT-4~\cite{openai2023gpt4} 的技术报告
        都使用了这种方法。
        局限是：它只能检测\textbf{精确重叠}，
        无法发现同义改写（paraphrase）形式的泄露。
        更大的局限是：它\textbf{需要访问训练数据}：
        对于 GPT-4 等闭源模型，外部研究者
        无法使用这种方法。
  \item \textbf{成员推断}（membership inference）：
        通过分析模型对特定样本的困惑度或置信度
        来判断该样本是否出现在训练数据中。
        直觉是：模型对训练数据中见过的文本
        会表现出异常低的困惑度。
  \item \textbf{CoDeC}（Contamination Detection via Context）：
        NVIDIA 在 2025 年提出的方法，
        通过比较\textbf{有无 in-context 示例}
        时模型对基准测试题目的置信度变化来检测污染。
        关键洞察是：对于未见过的数据，
        in-context 示例通常\textbf{提升}模型置信度；
        而对于已记忆的数据，
        in-context 示例反而可能\textbf{降低}置信度
        （因为干扰了记忆的检索模式）。
        % NEW REF: zawalski2025codec = Zawalski, Michal and Boubdir, Meriem and Balazy, Klaudia and Nushi, Besmira and Ribalta, Pablo. Detecting Data Contamination in LLMs via In-Context Learning. arXiv preprint arXiv:2510.27055, 2025.
  \item \textbf{动态基准测试}：
        与其检测污染，不如让污染无法发生：
        使用动态生成的基准测试题目（每次评估不同），
        使得训练数据中不可能包含未来生成的测试题。
        LiveBench 和 GAIA 等基准测试采用了这种策略。
  \item \textbf{Min-K\% Prob}（Shi 等人，2024）% NEW REF: shi2024minkprob = Detecting Pretraining Data from Large Language Models, Shi et al., ICLR 2024
        是一种不需要访问训练数据、
        也不需要参考模型的检测方法。
        它的核心洞察是：
        对于模型见过的文本，
        即使文本中最“意外”的 token
        （概率最低的 $K\%$ token），
        其对数概率也不会太低。
        相反，对于未见过的文本，
        最低概率 token 的对数概率显著更低。
        通过计算文本中最低 $K\%$ token 的平均对数概率，
        可以有效区分训练数据和非训练数据。
        在 GPT-2 和 LLaMA 上的实验中，
        Min-K\% Prob 的 AUROC 达到了 0.85--0.92：
        远优于传统的困惑度检测（AUROC 约 0.65）。
\end{itemize}

\subsection{实践中的去污染策略}

了解了检测方法后，
让我们看看前沿团队如何在实践中执行去污染。

LLaMA 3~\cite{dubey2024llama3} 使用了\textbf{三层去污染}流水线：
\begin{enumerate}[noitemsep]
  \item \textbf{8-gram 重叠去除}：
        对训练数据中与任何已知基准测试
        共享 $\geq$ 8 连续 token 的段落进行标记。
  \item \textbf{困惑度异常检测}：
        对训练数据在模型上的困惑度做分布分析，
        标记困惑度异常低的文档
        （可能是模型已经“记住”的内容）。
  \item \textbf{“保险”报告}：
        在所有基准测试结果中，
        同时报告“去污染前”和“去污染后”的分数，
        让读者自行判断污染的影响。
\end{enumerate}

GPT-4 的技术报告采用了更保守的策略：
使用 \textbf{substring match + fuzzy matching}
对训练数据进行大规模扫描，
并在报告中对“可能受污染影响”的基准测试结果
加注了详细的警告说明。
然而，由于 GPT-4 的训练数据未公开，
外部研究者无法独立验证其去污染的有效性：
这是\textbf{数据透明度}至关重要的又一个例子。

\textbf{2025--2026 年的新趋势}是
将去污染从“事后检测”前移到“预防性设计”。
具体做法包括：
\begin{itemize}[noitemsep]
  \item \textbf{训练数据索引化}：
        构建训练数据的高效索引
        （如 Bloom filter 或 MinHash 索引），
        使得在训练数据中查询特定文本
        只需毫秒级延迟。
        这让“实时去污染”成为可能：
        每当新的基准测试发布，
        可以立即检查并从训练数据中移除相关内容。
  \item \textbf{时间戳感知的训练数据管理}：
        为每条训练数据记录其来源的时间戳，
        然后规定“训练数据的时间截止线
        必须早于所有评估基准的发布日期”。
        这在理论上可以完全避免
        “训练数据包含基准测试数据”的问题：
        但前提是基准测试数据确实在
        截止日期之后才首次出现在互联网上。
\end{itemize}

\begin{warnbox}[数据污染是一个系统性问题]
数据污染不是个别团队的失误，
而是当前 LLM 训练范式的\textbf{系统性问题}。
当训练数据规模达到万亿 token 级别时，
要精确知道训练数据中包含什么几乎是不可能的。
即使团队真诚地尝试去除基准测试数据，
同义改写、翻译、引用等间接泄露渠道也难以完全堵截。
这使得\textbf{数据溯源}（data provenance）
和\textbf{基准测试设计}成为下一阶段 LLM 研究的重要课题。
\end{warnbox}

\section{版权与法律：数据的法律边界}

训练数据的法律问题已经从学术讨论
演变为影响整个 AI 行业的\textbf{核心挑战}。
2024--2026 年的一系列诉讼和立法
正在重塑 LLM 训练数据的获取方式。

\subsection{《纽约时报》诉 OpenAI 案}

2023 年 12 月，《纽约时报》对 OpenAI 和微软
提起了版权侵权诉讼：
这是迄今为止 AI 训练数据领域\textbf{影响最大}的案件。
《纽约时报》指控 OpenAI 在未经授权的情况下，
将其数百万篇受版权保护的新闻报道
用于训练 GPT 系列模型，
并展示了 ChatGPT 能够\textbf{逐字复述}
《纽约时报》文章内容的证据。

截至 2026 年初，此案仍在审理中。
2025 年 4 月的法庭裁决对《纽约时报》有利：
法官认为其部分诉求（包括指控 OpenAI 的训练构成侵权）
在现阶段具有成立的可能性。
这一裁决虽非最终判决，
但已经对行业产生了深远影响：
多家 AI 公司开始与内容出版商签订\textbf{许可协议}，
为训练数据付费。

\subsection{其他重要诉讼与判例}

《纽约时报》案并非孤例。
2024--2026 年间，AI 训练数据的版权诉讼已形成\textbf{诉讼浪潮}：

\begin{itemize}[noitemsep]
  \item \textbf{Getty Images 诉 Stability AI}（2023 年起）：
        图片库 Getty Images 起诉 Stable Diffusion 的开发商
        未经授权使用其受版权保护的图片训练模型。
        这一案件虽然主要涉及图像模型，
        但其判决对文本 LLM 同样具有\textbf{先例效应}：
        如果法院认定图像训练构成侵权，
        同样的逻辑将适用于文本训练。
  \item \textbf{作家联合诉讼}（Authors Guild v. OpenAI，2023 年起）：
        包括 John Grisham、Jonathan Franzen 等知名作家在内的
        数千名作家联合起诉 OpenAI，
        指控其使用 Books3 等未授权的书籍数据集训练模型。
        Books3 包含约 20 万本书的全文：
        其中绝大多数受版权保护。
        这一案件的核心争议是：
        The Pile 中的 Books3 子集来自盗版电子书网站 Bibliotik，
        即使 OpenAI 本身没有直接从盗版网站获取数据，
        使用由第三方“清洗”后的版本是否构成侵权？
  \item \textbf{Thomson Reuters 诉 Ross Intelligence}（2024 年裁决）：
        法院裁定 Ross Intelligence 使用 Westlaw 的法律文本训练 AI
        \textbf{不构成合理使用}：
        这是 AI 训练领域第一个对 AI 公司不利的正式裁决。
        虽然案件规模较小，
        但其“非合理使用”的判决结论对后续案件具有参考价值。
\end{itemize}

\subsection{合理使用（Fair Use）之争}

美国的核心法律争议集中在
\textbf{合理使用}（fair use）原则上。
合理使用是美国版权法第 107 条规定的抗辩理由，
法院需要权衡四个因素：
\begin{enumerate}[noitemsep]
  \item \textbf{使用的目的和性质}——是否为“转化性”使用？
  \item \textbf{原作品的性质}——是事实性还是创造性作品？
  \item \textbf{使用的数量}——使用了原作品的多大比例？
  \item \textbf{对市场的影响}——是否损害了原作品的市场价值？
\end{enumerate}

AI 公司在\textbf{因素 1}上有较强的论据：
LLM 训练确实是“转化性”的：
模型学习的是语言的统计模式，
而不是复制特定的文章。
类似的，Google 在 Google Books 案中
成功主张了对图书的扫描和索引构成合理使用。

但在\textbf{因素 4}上，内容创作者的论据同样有力：
LLM 训练与搜索引擎索引有本质区别：
搜索引擎引导用户访问原始内容（增加流量），
而 LLM 直接生成回答、替代原始内容（减少流量）。
如果用户可以直接向 ChatGPT 询问最新新闻，
为什么还要订阅《纽约时报》？
2025 年的市场数据似乎支持了这一担忧：
多家新闻机构报告，自 ChatGPT 发布以来，
来自搜索引擎的流量下降了 15--25\%，
部分归因于 AI 生成的摘要替代了用户的点击行为。

\textbf{因素 3}（使用数量）是一个特别微妙的维度。
OpenAI 主张，虽然他们使用了\textbf{整篇}文章训练，
但模型只学习了“统计模式”而非记忆了具体内容。
然而，《纽约时报》的律师展示了
ChatGPT 能够\textbf{逐字逐句复述}长达数百词的文章片段：
这强烈暗示模型\textbf{确实记住了}部分训练数据。
这一“逐字复述”的证据
对 OpenAI 的“仅学习统计模式”论据构成了严重挑战。

\subsection{数据溯源与 Opt-Out 机制}

法律压力正在催生一个全新的技术领域：
\textbf{数据溯源}（data provenance）。

\textbf{Opt-Out 机制}的演进是这一趋势的缩影。
传统的 \texttt{robots.txt} 是网站声明
“不希望被爬取”的唯一机制：
但它是一个\textbf{礼仪约定}而非法律强制。
2024 年后，随着欧盟 AI 法案的生效，
\texttt{robots.txt} 开始具有法律效力，
但其粗粒度（只能按整个网站或目录声明）
无法满足“允许搜索引擎索引但禁止 AI 训练”的需求。

为此，\textbf{TDM-reservation}（Text and Data Mining reservation）
标准被提出。
它通过在网页的 \texttt{<meta>} 标签中声明
“本页面不允许用于 AI/TDM 目的”，
提供了比 \texttt{robots.txt} 更精细的控制。
欧盟 AI 法案要求 AI 公司\textbf{必须尊重}
TDM-reservation 声明：
违反者可能面临重大罚款
（年全球收入的 3\%，对大型 AI 公司意味着数十亿美元）。

\textbf{C2PA}（Coalition for Content Provenance and Authenticity）
是另一个重要的标准。
C2PA 为数字内容提供了\textbf{不可篡改的来源记录}：
类似于区块链但应用于内容领域。
每段内容从创建到使用的整个生命周期
都被记录在 C2PA 的证书中。
Adobe、Microsoft、Google 等公司已经在
自家的产品中集成了 C2PA 支持：
例如，Photoshop 导出的图片会自动嵌入 C2PA 元数据，
记录其创建历史和任何 AI 编辑操作。
将 C2PA 扩展到 LLM 训练数据的溯源
是一个活跃的工程方向：
目标是让任何人都能查询
“这个模型的训练数据中包含了我的内容吗？”。

\subsection{欧盟 AI 法案与全球监管}

\textbf{欧盟 AI 法案}（EU AI Act，2024 年正式生效）
对 AI 训练数据提出了明确的\textbf{透明度要求}：
通用 AI 模型的提供者必须
“公布训练数据中受版权保护内容的充分详细摘要”。
虽然具体的执行标准仍在制定中，
但这一要求已经迫使 AI 公司开始构建
\textbf{数据溯源}（data provenance）系统：
记录训练数据中每一条内容的来源、许可证状态和使用限制。

欧盟还规定了\textbf{选择退出}（opt-out）机制：
版权持有者可以明确声明其内容不得用于 AI 训练
（例如通过 robots.txt 或 TDM-reservation 标签），
AI 公司必须尊重这些声明。
这使得 robots.txt 从一个非强制性的礼仪约定
变成了一个具有法律效力的 opt-out 机制。

\textbf{全球监管的碎片化}是另一个值得关注的趋势。
不同国家和地区对 AI 训练数据的法律态度差异巨大：

\begin{itemize}[noitemsep]
  \item \textbf{日本}：2018 年修订的《著作权法》
        明确规定“出于信息分析目的的复制”
        不构成侵权：
        这使得日本成为 AI 训练数据法律上
        \textbf{最友好}的主要经济体之一。
        但 2025 年日本政府开始考虑
        对这一宽松政策进行修订，
        特别是在生成式 AI 的输出
        与原作品“实质性相似”的情况下。
  \item \textbf{中国}：《生成式人工智能服务管理暂行办法》
        （2023 年 8 月生效）要求 AI 服务提供者
        “使用具有合法来源的数据和基础模型”。
        但“合法来源”的定义仍然模糊，
        实践中大量中文模型使用了
        未经明确授权的网页数据训练。
        2025 年的多起诉讼
        （如某网文平台诉 AI 公司使用网络小说训练模型）
        正在推动法律界限的明确化。
  \item \textbf{英国}：曾计划为 AI 训练引入广泛的版权例外，
        但在创意产业的强烈反对下于 2024 年搁置。
        目前英国的立场是“逐案判断”（case-by-case），
        没有为 AI 训练提供明确的法律安全港。
\end{itemize}

这种监管碎片化对\textbf{全球化的 AI 公司}
构成了巨大的合规挑战：
一套训练数据在日本完全合法、
在欧盟需要尊重 opt-out、
在美国取决于合理使用的判例、
在中国需要“合法来源”。
这也解释了为什么越来越多的公司
选择构建\textbf{区域差异化}的训练数据管线：
为不同市场准备不同的数据组合，
而非使用单一的全球数据集。

\subsection{行业的应对策略}

面对法律风险，AI 行业采取了多种应对策略：

\begin{itemize}[noitemsep]
  \item \textbf{许可协议}：
        OpenAI 已与美联社、Axel Springer、Le Monde 等
        多家媒体机构签订了内容许可协议。
        Google 和 Anthropic 也在推进类似的合作。
  \item \textbf{合成数据替代}：
        用合成数据替代受版权保护的真实数据：
        这是 Phi 系列和 Nemotron-4 策略的另一个动机。
  \item \textbf{数据溯源基础设施}：
        构建能够追踪每条训练数据来源和许可证状态的系统。
        C2PA（Coalition for Content Provenance and Authenticity）
        等标准正在被用于标记 AI 生成的内容
        和追踪内容的使用历史。
  \item \textbf{开放数据运动}：
        社区推动创建明确授权用于 AI 训练的开放数据集：
        如 The Stack（仅包含宽松许可证的代码）、
        Dolma（完全透明的数据来源记录）。
\end{itemize}

\begin{corebox}[版权问题的长期影响]
版权争议可能从根本上改变 LLM 的训练范式。
如果法院最终裁定未经授权使用版权内容构成侵权，
那么 AI 公司将面临两个选择：
(1) 为训练数据支付许可费用：
这将大幅增加训练成本，
可能使得只有资金雄厚的公司才能训练 frontier 模型；
(2) 完全依赖公共领域和合成数据：
这将加速合成数据技术的发展，
但也增加了模型坍缩的风险。
无论哪种结果，
\textbf{数据溯源和透明度}
都将成为未来 LLM 训练的必备基础设施。
\end{corebox}

\section{分词：从文本到 Token}

数据处理流水线的最后一步是\textbf{分词}（tokenization）：
将清洗后的文本转换为模型能理解的 token 序列。
这一步虽然经常被忽视，但分词方案的选择
对模型的训练效率和最终能力有深远影响。

\subsection{BPE 分词}

\textbf{Byte Pair Encoding}（BPE）
是几乎所有现代 LLM 使用的分词方法。
BPE 的核心思想是\textbf{贪心合并}：
从单字节出发，
反复合并出现频率最高的相邻字节对，
直到词表达到目标大小。

BPE 的训练过程：
\begin{enumerate}[noitemsep]
  \item 初始化：词表为所有 256 个字节。
  \item 统计所有相邻 token 对的出现频率。
  \item 合并频率最高的一对，创建新 token。
  \item 重复步骤 2--3，直到词表达到目标大小
        （如 128K tokens）。
\end{enumerate}

词表大小是一个关键的设计选择。
更大的词表意味着更多的常见词或短语
被编码为单个 token：
相同的文本需要更少的 token 表示，
训练和推理更高效。
但更大的词表也意味着更大的嵌入矩阵
（$V \times d$，$V$ 为词表大小，$d$ 为嵌入维度），
这在小模型中占据参数量的很大比例。

主流模型的词表大小差异不小：

\begin{center}\small
\renewcommand{\arraystretch}{1.3}
\begin{tabular}{lrl}
\toprule
\rowcolor{figLightGray}
\textbf{模型} & \textbf{词表大小} & \textbf{分词器} \\
\midrule
GPT-2 & 50,257 & BPE (byte-level) \\
LLaMA 1/2 & 32,000 & SentencePiece \\
LLaMA 3 & 128,000 & tiktoken (BPE) \\
DeepSeek-V3 & 128,000 & BPE \\
Qwen3 & 151,936 & BPE \\
\bottomrule
\end{tabular}
\end{center}

LLaMA 从 32K 到 128K 的跳跃值得注意。
LLaMA 1/2 的 32K 词表对\textbf{非拉丁语言}不友好：
一个中文字符可能被分为 2--3 个 token，
而一个英文单词通常只需 1--2 个 token。
这意味着处理中文文本时，
模型“看到”的有效上下文长度只有英文的一半。
LLaMA 3 扩大到 128K 词表后，
中文的 token 效率提升了约 50\%：
同一段中文文本需要的 token 数减少了三分之一。
这对多语言能力的提升有直接贡献。

\subsection{多语言分词的挑战}

多语言分词面临\textbf{token 公平性}的问题。
在英文为主的训练语料上训练的 BPE 分词器
会优先为英文常见词和词组创建 token，
导致英文的 token 效率远高于其他语言。
Petrov 等人（2024）的研究% NEW REF: petrov2024tokenizer = Language Model is Not for All: A Survey of Tokenization in Large Language Models, Petrov et al., 2024
系统地量化了这一不公平性：
将同一篇 United Nations 文件翻译成不同语言后，
英文版需要约 100 个 token，
中文版需要约 150 个 token，
缅甸文版需要约 \textbf{450} 个 token：
4.5 倍的效率差距意味着缅甸语用户
在相同的上下文窗口中能输入的内容
不到英文用户的四分之一。

解决这一问题的常见策略是
在\textbf{多语言均衡语料}上训练分词器。
Qwen3 的 151K 词表就是在
精心均衡的多语言语料上训练的：
确保中文、日文、韩文等高需求语言
获得足够的 token 分配。
这是 Qwen 系列模型在中文任务上
表现优于 LLaMA 的原因之一：
不仅训练数据更多中文内容，
分词器本身也对中文更友好。

\subsection{序列打包}

分词之后的最后一步是\textbf{序列打包}（sequence packing）。
训练效率要求每个 batch 中的序列长度一致
（以充分利用 GPU 的并行计算能力），
但自然文档的长度差异很大：
从几十个 token 到数万个 token 不等。

两种主要的打包策略：
\begin{itemize}[noitemsep]
  \item \textbf{截断 + 填充}（truncation + padding）：
        将长文档截断到最大序列长度，
        短文档填充 \texttt{[PAD]} token 到最大长度。
        简单但浪费——填充的 token
        不贡献任何训练信号，
        却占据了计算资源。
  \item \textbf{多文档拼接}（document concatenation）：
        将多个短文档拼接成一个长序列，
        用特殊的 \texttt{<|endoftext|>} token 分隔。
        不浪费计算，但引入了跨文档的注意力：
        序列末尾的 token 可能“看到”前一个文档的内容，
        即使两者在语义上毫无关系。
\end{itemize}

几乎所有现代 LLM 训练都使用多文档拼接，
因为填充造成的计算浪费在大规模训练中不可接受。
跨文档注意力的问题通过
\textbf{注意力掩码}或 \textbf{document boundary markers}
来缓解——确保每个文档只与自身产生注意力。
LLaMA 3 的实现中，使用了
“cross-document attention masking”来隔离拼接的文档，
实验表明这比不做隔离在下游任务上
带来了约 0.5\% 的一致性提升。

\begin{corebox}[数据工程的核心原则]
\begin{enumerate}[noitemsep]
  \item \textbf{质量 $>$ 数量}：1T 高质量 token $>$ 10T 低质量 token。
        Phi 系列和 FineWeb-Edu 反复证明了这一点。
  \item \textbf{多样性至关重要}：单一来源的数据（即使高质量）会限制模型的泛化能力。
        混合网页、代码、学术、书籍等多种来源。
  \item \textbf{代码是推理能力的催化剂}：即使你的目标不是代码生成，也应该纳入代码数据。
        代码的逻辑严密性能提升模型的通用推理能力。
  \item \textbf{去重是必须的}：重复数据不仅浪费计算，还会降低模型的泛化能力。
        多层次去重（URL $\to$ 精确 $\to$ 模糊 $\to$ 子串）效果最好。
  \item \textbf{合成数据是补充而非替代}：始终混入真实数据以避免模型坍缩。
        最优比例约为 30\% 合成 + 70\% 真实。
  \item \textbf{可复现性是底线}：记录每一步的参数和随机种子，
        确保同事或后来者能精确复现你的数据流水线。
  \item \textbf{法律合规不可忽视}：了解训练数据的版权状态，
        构建数据溯源系统，
        尊重 opt-out 声明：
        法律风险正在成为数据工程的核心约束之一。
\end{enumerate}
\end{corebox}


%% ============================================================
%% --- NEW REFERENCES (expansion update) ---
% NEW REF: dohmatob2024collapse = Dohmatob, Elvis and others. A Tale of Tails: Model Collapse as a Change of Scaling Laws. arXiv preprint arXiv:2402.01942, 2024.
% NEW REF: guo2024collapse = Guo, Rui and others. Curious Decline of Linguistic Diversity: A Large-Scale Analysis of LLM-Generated Text. arXiv preprint arXiv:2404.10040, 2024.
% NEW REF: chen2024datajuicer = Chen, Daoyuan and others. Data-Juicer: A One-Stop Data Processing System for Large Language Models. In SIGMOD, 2024.
